% HELUT living textbook --- canonical TeX.
% Build: `make textbook` → textbook/helut-living-textbook.pdf + .md
% Markdown is generated; edit this tree only.
% Living epoch and claim IDs: directives/claim-sheet.md
%
% Chapters are \input (not \include) so latexmk -outdir stays well-behaved.

\documentclass[11pt,oneside,openany]{book}

% Shared preamble for the HELUT living textbook.
% Canonical TeX. Build with: make textbook

\usepackage[margin=1in]{geometry}
\usepackage{amsmath,amssymb,amsthm}
\usepackage{booktabs}
\usepackage{longtable}
\usepackage{graphicx}
\usepackage{hyperref}
\usepackage{microtype}
\usepackage{enumitem}
\usepackage{xcolor}
\usepackage{listings}
\usepackage{url}
\usepackage[T1]{fontenc}
\usepackage{textcomp}
\usepackage{fancyhdr}

% ---------------------------------------------------------------------------
% Draft status is carried structurally, not by hedging the prose.
%   1. \draftnotice  -- one unmissable page in the front matter.
%   2. Per-page footer stamp -- so no single page, quoted or printed alone,
%      can be mistaken for a finished textbook.
% The body is written in the voice of the finished book on purpose. Editorial
% readiness and claim correctness are different things: claims are still gated
% by directives/claim-sheet.md and Scripts/claim_audit.py.
% ---------------------------------------------------------------------------
\newcommand{\draftstamp}{%
  \footnotesize\itshape Draft \corpusedition{} (\livingepoch{}) --- in
  preparation%
}
\pagestyle{fancy}
\fancyhf{}
\fancyfoot[C]{\thepage}
\fancyfoot[L]{\draftstamp}
\renewcommand{\headrulewidth}{0pt}
\renewcommand{\footrulewidth}{0.4pt}
\fancypagestyle{plain}{%
  \fancyhf{}%
  \fancyfoot[C]{\thepage}%
  \fancyfoot[L]{\draftstamp}%
  \renewcommand{\headrulewidth}{0pt}%
  \renewcommand{\footrulewidth}{0.4pt}%
}

\hypersetup{
  colorlinks=true,
  linkcolor=black,
  citecolor=black,
  urlcolor=blue!50!black,
  pdftitle={Reconfigurable Homomorphic Computing: A Living Textbook},
  pdfauthor={Digital Defiance / HELUT}
}

\lstset{
  basicstyle=\ttfamily\small,
  breaklines=true,
  columns=fullflexible,
  frame=single,
  rulecolor=\color{black!20},
  showstringspaces=false
}

% ---------------------------------------------------------------------------
% Living epoch --- bump when claim-sheet.md gains a C-row or closes an H.
% Canonical inventory: directives/claim-sheet.md
% ---------------------------------------------------------------------------
\newcommand{\livingepoch}{2026-08-14 / C69}
\newcommand{\corpusedition}{0.1.2}
% Optional audit stamp from Scripts/helut_grand_audit.sh
\InputIfFileExists{git-stamp.tex}{}{%
  \newcommand{\gitstamp}{dirty}
  \newcommand{\auditepoch}{\livingepoch}
  \newcommand{\audittime}{unaudited}
}

% The corpus names its theorems "Theorem 1", "Theorem 1'", "Theorem 2", ...
% amsthm numbers them by chapter ("Theorem 9.1"). Both names are load-bearing:
% the corpus name is what directives/, the claim sheet and the papers cite, and
% the chapter number is what a reader of this book sees. \corpusthm bridges
% them inside the theorem's title so neither name is orphaned.
\newcommand{\corpusthm}[1]{corpus Theorem~#1}

% Claim-sheet IDs
\newcommand{\cid}[1]{\textbf{C#1}}
\newcommand{\hid}[1]{\textbf{H#1}}
\newcommand{\nid}[1]{\textbf{N#1}}

% Pointers into the HELUT corpus (repo-relative)
\newcommand{\corpus}[1]{\texttt{#1}}

% Status boxes: reproduced / hedge / frontier / grows-when
\newcommand{\livingbox}[3]{%
  \par\medskip
  \noindent\fcolorbox{#1}{white}{%
    \begin{minipage}{\dimexpr\linewidth-2\fboxsep-2\fboxrule\relax}
    \smallskip
    \textbf{#2.} #3
    \smallskip
    \end{minipage}%
  }\par\medskip
}
\newcommand{\reproduced}[2]{\livingbox{green!50!black}{Reproduced (#1)}{#2}}
\newcommand{\hedgebox}[2]{\livingbox{orange!70!black}{Open hedge (#1)}{#2}}
% Measured negative: a real result that did NOT pass. Doctrine prints clean
% negatives, but they must never wear the green Reproduced label.
\newcommand{\negativebox}[2]{\livingbox{purple!60!black}{Measured negative (#1)}{#2}}
% Independent re-check by a second model. Weaker than a human gate; stronger
% than nothing. Never a substitute for \reproduced.
\newcommand{\recheckbox}[2]{\livingbox{teal!60!black}{Re-checked, second model (#1)}{#2}}
% Generic claim-kind noun, so \cid{} is never called with an empty argument
% and stays machine-extractable.
\newcommand{\claimkind}[1]{\textbf{#1}}
\newcommand{\frontierbox}[1]{\livingbox{blue!55!black}{Frontier --- not a claim}{#1}}
\newcommand{\growswhen}[1]{\livingbox{black!45}{This section grows when}{#1}}
\newcommand{\nonclaimbox}[1]{\livingbox{red!55!black}{Non-claim}{#1}}
\newcommand{\scaffoldbox}[1]{\livingbox{red!70!black}{Not a course yet}{#1}}

% The single authoritative statement of editorial status. Everything that reads
% like a finished course later in the book is covered by this page.
\newcommand{\draftnotice}{%
  \thispagestyle{plain}%
  \begin{center}
  \fcolorbox{red!70!black}{red!4}{%
    \begin{minipage}{\dimexpr\textwidth-4\fboxsep-4\fboxrule\relax}
    \vspace{0.6em}
    \begin{center}\Large\bfseries Draft in preparation --- not yet ready to teach from\end{center}
    \vspace{0.2em}
    \normalsize
    Edition \corpusedition{}, epoch \livingepoch{}.

    \medskip
    This volume is being \emph{prepared} as a textbook, and is written in the
    voice of the book it is intended to become. It therefore may contain text
    that implies a course already exists --- chapters that lecture, exercises
    that are set, labs that are specified, a syllabus that assigns weeks and
    grading weights. Please read that as work in progress. The material has not
    been taught, has not been reviewed by anyone qualified to sign off on the
    pedagogy, and is not part of any curriculum.

    \medskip
    Status is stated here once and stamped on each page rather than repeated
    inside every paragraph, so that the text can be read --- and grown --- as
    continuous prose.

    \medskip
    \textbf{What is nonetheless load-bearing.} Editorial readiness and factual
    correctness are separate. Every \cid{}, \hid{} and \nid{} reference, every
    \textbf{Reproduced} box, and every theorem in this book is bound to
    \corpus{directives/claim-sheet.md} and a reproduce receipt, and that binding
    is checked mechanically by \corpus{Scripts/claim\_audit.py} in CI. Claims,
    lemmas, theorems and formulae are not invented to fill the scaffold: a human
    checks each against the sheet before commit. If a sentence in this book
    disagrees with the claim sheet, \textbf{the sheet wins} and the sentence is
    a bug.

    \medskip
    Numbers carry the same discipline. Timings are medians of repeated runs with
    the spread printed where they have been re-measured
    (\corpus{Scripts/bench\_repeat.py}); single-observation figures inherited
    from earlier work are being replaced incrementally.
    \vspace{0.6em}
    \end{minipage}%
  }
  \end{center}
  \clearpage
}

% Pedagogy.
%
% These are real exercises and real labs, written in the voice the finished book
% will use. That is deliberate: this edition is a DRAFT being prepared toward a
% course, not a course pretending to be a draft. Editorial status is carried
% once, loudly, by the front-matter notice and the per-page DRAFT stamp -- not
% by hedging every environment title. See \draftnotice below.
\theoremstyle{definition}
\newtheorem{definition}{Definition}[chapter]
\newtheorem{example}{Example}[chapter]
\newtheorem{exercise}{Exercise}[chapter]
\newtheorem{lab}{Laboratory}[chapter]
\theoremstyle{plain}
\newtheorem{lemma}{Lemma}[chapter]
\newtheorem{theorem}{Theorem}[chapter]
\newtheorem{proposition}{Proposition}[chapter]
\theoremstyle{remark}
\newtheorem{remark}{Remark}[chapter]
\newtheorem{warning}{Honesty note}[chapter]

\setlist[itemize]{leftmargin=*}
\setlist[enumerate]{leftmargin=*}


\title{\textbf{Reconfigurable Homomorphic Computing}\\
\large A Living Textbook of Netlist-Clocked FHE,\\
Differentiable Hardware, and Polymorphic Ciphers}
\author{Digital Defiance\\
\textit{HELUT Project --- living edition \corpusedition{} (\livingepoch)}}
\date{August 2026}

\begin{document}
\frontmatter

\begin{titlepage}
\thispagestyle{empty}
\centering
\vspace*{2cm}
{\Huge\bfseries Reconfigurable Homomorphic Computing\par}
\vspace{1.2em}
{\Large A Living Textbook\par}
\vspace{0.8em}
{\large Netlist-Clocked Torus FHE \textperiodcentered{}
Differentiable Hardware \textperiodcentered{}
Adversarial Polymorphic Ciphers\par}
\vfill
{\large Digital Defiance / HELUT Project\par}
\vspace{0.5em}
{\large Living edition \corpusedition{} \textperiodcentered{} epoch \livingepoch\par}
\vspace{0.2em}
{\small Audit stamp: \texttt{\gitstamp} \textperiodcentered{} \audittime\par}
\vspace{0.3em}
{\large Scaffold for a future course --- not ready to teach\par}
\vfill
{\small Canonical source: \texttt{textbook/helut-living-textbook.tex}.
Corpus of record: \texttt{directives/claim-sheet.md}.
This volume is an AI-bootstrapped collecting place for theorems and
formulae. Humans must vet, bulk, fill, and edit it before it is a
textbook. Frontier chapters are labeled frontier.\par}
\end{titlepage}

\draftnotice

\chapter*{How to read the living boxes}
\addcontentsline{toc}{chapter}{How to read the living boxes}

Seven boxes appear throughout this book.
They are not decoration; they are the difference between a textbook and a press release.
Editorial status is not their job --- that is stated once on the previous page and
stamped on every page footer. These boxes carry \emph{evidential} status, which is
a separate axis and the one that must never slip.

\reproduced{C\textit{n}}{A result that someone else can re-run from \corpus{REPRODUCE.md} and the claim inventory. The prose may treat it as established \emph{in this corpus}. Keyed to a \claimkind{C} row, always --- never to free text and never to an \claimkind{H} id alone.}

\hedgebox{H\textit{n}}{An honest asterisk. Say it with the hedge attached. Do not let it get dropped in a problem set or a paper.}

\negativebox{no \claimkind{C} row}{A measured \emph{failure}: something we ran that did not pass. Printing clean negatives is doctrine, so these appear on purpose --- but they are never green, and they are never receipts.}

\recheckbox{second model}{Independently re-derived by a different model reading the same evidence. Weaker than a human check and stronger than nothing. Never a substitute for \textbf{Reproduced}.}

\frontierbox{Speculative territory from \corpus{directives/potential-avenues.md}. A seminar essay, never a midterm fact.}

\growswhen{A stub waiting on a new \claimkind{C} row, a closed \claimkind{H}, or a graduated avenue. Skip, or read the matching directive.}

\scaffoldbox{Section-level reminder that a chapter is skeleton rather than
finished text. The book-level statement of that is the draft notice on the
previous page; this box marks the specific places where the scaffolding shows.}

\paragraph{Why theorems carry two numbers.}
The corpus --- \corpus{directives/}, the claim sheet, the papers --- has called
these results \emph{Theorem 1}, \emph{Theorem 1$''$}, \emph{Theorem 2} and
\emph{Theorem 3} since before this book existed, and external notes cite them
that way. \LaTeX{} numbers theorems by chapter. Rather than break either
convention, each theorem prints both: Theorem~9.1 of this book is also
\corpusthm{1}. If you are following a reference from the claim sheet, match on
the corpus name; if you are reading straight through, use the chapter number.

\tableofcontents

\mainmatter

\part{How This Book Lives}

\chapter{Preface}
\label{ch:preface}

This book is the first teaching surface for a subject that did not have a name
until the laboratory forced one.

Universities already teach \emph{reconfigurable computing}: LUTs, interconnect,
place-and-route, FPGA fabrics, high-level synthesis.
They already teach \emph{applied cryptography}: AES, RSA, lattices, sometimes
a lecture on fully homomorphic encryption.
They almost never teach what happens when the reconfigurable fabric \emph{is}
a machine-learning graph API, the bitstream is a Yosys netlist, and the values
on the wires are torus samples that must remain encrypted while a host clock
ticks flip-flops.

That intersection is the subject of this course.
We call the working stack \textbf{HELUT}---Homomorphic Edge Look-Up Tensors---and
we call the subject \textbf{reconfigurable homomorphic computing}.

The stack has three pillars~\cite{helut-release,helut-paper}:

\begin{enumerate}
  \item \textbf{Netlist-clocked torus FHE.}
        Gate-level sequential circuits evaluate as exact
        $\mathbb{Z}/2^{32}\mathbb{Z}$ tensor graphs, with real bootstrapping-key
        blind rotation and machine-checkable certificates.
  \item \textbf{Differentiable hardware cryptanalysis.}
        Discrete INIT tables melt into continuous tensors; a $\lambda$-squeeze
        grades shatter versus hold; reciprocal structure that combinatoric
        sieves miss can reappear as an involution.
  \item \textbf{Adversarial polymorphic ciphers.}
        Blue ciphers co-evolve under Red pressure and fail closed, rather than
        waiting to become the next historically leaked stream cipher.
\end{enumerate}

\paragraph{Why a living textbook.}
Research papers freeze a slice of a laboratory.
This course cannot.
The claim inventory~\cite{helut-claim-sheet} moves: Metal kernels get faster,
hedges close, avenues stay unlabeled until they earn receipts.
A professor who taught from a PDF dated June would be teaching a different
Metal compiler than a professor teaching from August.
So the book is versioned against the corpus \emph{epoch}
(\livingepoch{} in this edition) and is honest about stubs.

\paragraph{What this edition is.}
Edition~\corpusedition{} is the \emph{beginnings}: enough foundations to open a
course, enough Pillar~I to run laboratories, enough Pillars~II and~III to
assign readings, and a frontier part that is labeled as frontier.
It is not a complete treatise.
Where the laboratory is thin, the page says so.

\paragraph{What this edition is not.}
It is not a claim that HELUT invented TFHE~\cite{chillotti2020tfhe,ducas2015fhew}.
It is not a production key-size manual until a lattice estimator fills the
pending table (\hid{1}).
It is not a decrypt of U-534 / P1030680.
It is not a side-channel paper.
Those sentences are printed again in Chapter~\ref{ch:living} because students
will forget them the week before the final.

\paragraph{Who wrote it.}
The HELUT project at Digital Defiance.
The scientific voice of the book is the same voice as the research-release
doctrine: lead with what runs, and with what it does not prove.

\paragraph{How to cite an epoch.}
Cite the git commit and the claim-sheet epoch, not ``the HELUT textbook'' as if
it were a static ISBN.
When a student reproduces \cid{18}, they are reproducing a dated Metal NTT
persist tile, not an eternal constant of nature.

\chapter{To the instructor}
\label{ch:instructor}

You are being asked to teach a class that does not yet appear in the ACM/IEEE
curriculum guidelines under that name.
This chapter is the briefing.

\section{The course in one paragraph}

Students learn to treat a \emph{synthesized gate netlist} as the universal
intermediate representation of reconfigurable computing, to lower that netlist
onto a \emph{tensor engine} (here: Apple Silicon \texttt{MPSGraph}), and to
evaluate it either as a boolean-faithful oracle or as torus FHE with
certificates.
They then learn a second compiler direction---melting discrete hardware into
continuous tensors---and a third: designing ciphers that are born under
adversarial pressure.
The moral of the course is not ``FHE is fast on a laptop.''
The moral is: \emph{reconfigurable computing has a new fabric, and honesty about
what that fabric proves is part of the engineering.}

\section{Prerequisites}

\begin{center}
\begin{tabular}{@{}lp{0.62\textwidth}@{}}
\toprule
\textbf{Required} & Digital logic (LUTs, flip-flops, timing)~\cite{harris2012digital};
linear algebra; a first crypto course (modular arithmetic, attack models)~\cite{paar2009crypto}. \\
\textbf{Helpful} & One HDL (Verilog or VHDL); one systems language; Yosys as a
black box~\cite{yosys}. \\
\textbf{Not required} & Prior FHE, prior Metal, prior Swift, prior lattice
cryptography. The torus chapter is self-contained at the level this course needs.
\\
\bottomrule
\end{tabular}
\end{center}

\section{Two tracks}

Not every department can put an M-series Mac in front of every student.

\begin{description}
  \item[Systems track.]
        Apple Silicon laboratory.
        Students build, tick, and certificate-read HELUT from
        \corpus{REPRODUCE.md}.
        Labs~\ref{ch:lab} assume this track.
  \item[Theory track.]
        No GPU required.
        Students work the five-cell test, noise lemmas, and published logs as
        \emph{data}.
        A department can still grade a claim-audit final (below).
\end{description}

A mixed offering works: theory problem sets for everyone; systems labs as a
project option or a second section.

\section{Suggested grading}

\begin{center}
\begin{tabular}{@{}lp{0.22\textwidth}p{0.48\textwidth}@{}}
\toprule
Weight & Instrument & Why it exists \\
\midrule
30\% & Problem sets & Negacyclic arithmetic, packing, certificate reading. \\
40\% & Laboratories or theory equivalents & Receipts, not vibes. \\
20\% & \textbf{Claim audit} & Take a public sentence; run the five-cell test
from Chapter~\ref{ch:living}. Fail the assignment if the student upgrades a
hedge to a theorem. \\
10\% & Frontier essay & Pick one avenue from Chapter~\ref{ch:frontier}.
Must label it frontier. Must propose the \emph{first experiment} that would
graduate it. \\
\bottomrule
\end{tabular}
\end{center}

The claim audit is the signature assignment of this course.
It trains the skill the field currently lacks: refusing to launder an
open hedge into a keynote.

\section{What to lecture, what to assign as reading}

\begin{itemize}
  \item \textbf{Lecture} Parts~II--III (foundations and Pillar~I), the
        honesty chapter, and Theorem~\ref{thm:tensorlut} (\cid{19}).
  \item \textbf{Seminar} Pillar~II empirics (shatter / hold, campaign as
        case study) and Pillar~III in edition~\corpusedition{}.
        Do not lecture shatter grades as if they were Theorem~\ref{thm:tensorlut}.
  \item \textbf{Do not lecture as fact} Chapter~\ref{ch:frontier}.
        Give it as the last two seminars, with the frontier boxes left on
        the slides.
\end{itemize}

\section{How to pull a new edition mid-semester}

Science will move under your feet.
That is the point of a living textbook.

\begin{enumerate}
  \item Diff \corpus{directives/claim-sheet.md} against the epoch printed on
        the title page.
  \item If a new \cid{} row appeared, assign the matching lab receipt as a
        one-page addendum; do not pretend the old PDF contained it.
  \item If an \hid{} closed, tell the class explicitly: ``this asterisk is
        gone; here is the log.''
  \item If only frontier text moved, do not change the midterm.
\end{enumerate}

Bump \texttt{\textbackslash livingepoch} in \corpus{textbook/preamble.tex}
when you accept an addendum.
Cite git.

\section{Classroom norms}

\begin{itemize}
  \item Trivial / mock torus encodings are a \emph{shape laboratory}, not FHE.
        Students who call the boolean oracle ``encrypted computing'' lose
        points on purpose.
  \item \textbf{Campaign cryptanalysis} (Enigma / P1030680) is a \emph{case study} in
        Pillar~II methods, not the FHE claim, and not a decrypt.
  \item \textbf{Do not lecture} HELUT as an unsupervised or mobile agent, a
        self-rewriting FPGA, or a production encrypted CPU. Those sentences fail
        the five-cell test of Chapter~\ref{ch:living}.
  \item \hid{4} Grade~B is covering Track~A (\cid{52}--\cid{54}), not native
        $k{=}1$, not \texttt{cryptoPublicMS}, not PicoRV covering at production
        $N$.
  \item Speculative systems (encrypted LLM firewalls, self-modifying dark
        binaries, honey-token ledgers) are ethics-and-design discussions,
        not implementation homework, until they have \cid{} rows.
\end{itemize}

\growswhen{A published instructor's guide with slide skeletons, a Moodle/Canvas
quiz bank, and a non-Apple reference implementation of the negacyclic oracle
for the theory track.}

\chapter{The living contract}
\label{ch:living}

\section{Why a textbook needs a contract}

Ordinary textbooks accrete errata.
This one accretes \emph{science}.
Without a contract, the lecture slides of 2027 will still be quoting a fused
\texttt{MPSGraph} megagraph that the laboratory killed in August 2026, or---worse---will
quote a speculative avenue as if it were a theorem.

The contract is the research-release doctrine~\cite{helut-release} translated
into pedagogy.

\section{The three inventories}

Every public sentence in HELUT is classified in
\corpus{directives/claim-sheet.md}:

\begin{center}
\begin{tabular}{@{}llp{0.52\textwidth}@{}}
\toprule
Tag & Name & Classroom meaning \\
\midrule
\cid{$n$} & Reproducible result & You may lecture it. Students may cite it.
Someone else can re-run it~\cite{helut-claim-sheet}. \\
\hid{$n$} & Open hedge & You may discuss it. You must keep the asterisk.
A problem-set ``true/false'' that drops the asterisk is false. \\
\nid{$n$} & Non-implication & You must actively prevent this inference.
It is a common exam trap. \\
\bottomrule
\end{tabular}
\end{center}

A fourth class lives only in this book and in
\corpus{directives/potential-avenues.md}:

\begin{center}
\begin{tabular}{@{}lp{0.62\textwidth}@{}}
\toprule
Frontier & Not a claim until it has a \cid{} row and a reproduce command.
Seminar fuel. Never a midterm fact. \\
\bottomrule
\end{tabular}
\end{center}

\section{The five-cell test}

For an atomic claim $C$, the doctrine requires all five, or $C$ is hedged:

\begin{enumerate}
  \item \textbf{Proof} --- lemma, reduction, or graded empirical protocol with
        pass/fail.
  \item \textbf{Table} --- parameters, comparisons, or grades.
  \item \textbf{Metric} --- a number with units and a bar.
  \item \textbf{Examples} --- at least two concrete, runnable instances.
  \item \textbf{Application} --- at least one end-to-end reason to care.
\end{enumerate}

\begin{example}[Five cells for \cid{18}]
Metal 3-prime NTT persist blind-rotate.
Proof: CPU$+$Metal NTT $\equiv$ schoolbook (19 tests) plus microbench PASS bits.
Table: $N\in\{64,1024\}$ s/BR versus schoolbook persist and fused-EP.
Metric: $0.433\,\mathrm{s}$/BR at $N=1024$ (gpu $0.43\,\mathrm{s}$, RSS $148\,\mathrm{MiB}$).
Examples: microbench identity LUT; full\_adder boolean
SING (\textbf{S}timuli \textbf{In}, \textbf{G}olden out---encrypted evaluation
must match the cleartext netlist on the same generic vectors; see
Chapter~\ref{ch:pillari}).
Application: production-shaped polynomial degree on a laptop GPU graph.
\end{example}

\begin{example}[Five cells for \cid{19}]
TensorLUT continuous$\to$discrete Theorem~1~\cite{helut-tensorlut-thm}.
Proof: six \texttt{TensorLUTFormal.check*} lemmas, aggregated by
\texttt{certificate()}.
Table: lemma vs \texttt{holds}; campaign grades kept separate.
Metric: $F_{\mathrm{crypto}}=0$ on unmutated baseline; $\pi=0$ on binary INIT.
Examples: M4 baseline emit; freeze-core involution; blind 3-pair PASS.
Application: a stecker search that cannot propose a non-reciprocal map.
Asterisk: structural, not a U-534 break (\hid{6}).
\end{example}

\begin{exercise}[Claim audit]
Take any sentence from a blog post, a grant abstract, or last year's slide
deck about ``encrypted computing on a Neural Engine.''
Fill the five cells or mark the sentence a non-claim.
If you cannot name the log file, it is not a \cid{} row.
\end{exercise}

\section{Disclosure versus discovery}

\begin{center}
\begin{tabular}{@{}lp{0.38\textwidth}p{0.38\textwidth}@{}}
\toprule
 & Disclosure package & Trajectory \\
\midrule
Goal & Someone else re-runs what we assert & We know what to invent next \\
Artifact & \corpus{REPRODUCE.md}, logs, reproduced boxes
  & trajectory file, Chapter~\ref{ch:open} \\
Failure mode & Unreproducible hype & Silence / no next experiment \\
\bottomrule
\end{tabular}
\end{center}

When a trajectory item graduates, it earns a \cid{} row (or a tightened
\hid{}) and a reproduce command---\emph{then} it may enter lecture voice.
That sentence is the entire editorial policy of this book.

\section{Standing non-claims}

Print these on the course website.

\nonclaimbox{HELUT does not introduce a new lattice assumption, does not claim
to have invented TFHE, and does not offer ``176-bit secure'' as a slide title.
\cid{23} filled the Sage JSON (production HELUT $175.7$ vs Sage $180.2$ on
\texttt{prod-n1024-s16}); \hid{1} is the remaining $|\Delta|>16$ anchors and the
core-SVP vs Cost \texttt{rop} split.
Trivial/oracle Metal graphs are not FHE.
Covering Track~A Grade~B (\cid{52}--\cid{54}) is not native $k{=}1$ and
not \texttt{cryptoPublicMS}.
\texttt{publicMS} gadgets are not a side-channel theorem.
Campaign fitness is not encrypted tick rate.
TensorLUT involution grades are not a U-534 break.
Quantum attacks are out of scope.
Metal/GPU power analysis is parked~\cite{helut-avenues}.}

\section{How a chapter grows}

\begin{enumerate}
  \item A laboratory result lands in a log under \corpus{logs/}.
  \item A \cid{} or \hid{} row is updated in the claim sheet.
  \item The matching chapter gains an example, a table row, or a converted
        box; \texttt{\textbackslash livingepoch} bumps.
  \item The syllabus appendix does \emph{not} automatically change.
        Instructors opt in (Chapter~\ref{ch:instructor}).
\end{enumerate}

\growswhen{Automated extraction of claim-sheet tables into
Appendix~\ref{app:claims}, so the snapshot cannot drift.
Until then, humans diff the sheet against the appendix before each teaching
epoch.}


\part{Foundations of the Subject}

\chapter{A new subject}
\label{ch:subject}

\section{The gap between two mature fields}

Reconfigurable computing classically means: a designer writes RTL, a synthesizer
emits LUTs and flip-flops, a bitstream configures an FPGA, and the fabric
evaluates the circuit at a clock~\cite{harris2012digital,gajski2009embedded}.
The fabric is spatially wired. The configuration is a bitstream. The values
on the wires are bits.

Fully homomorphic encryption classically means: a scheme evaluates circuits
\emph{on ciphertexts}~\cite{chillotti2020tfhe,ducas2015fhew}.
The ``fabric'' is an arithmetic library---NTT butterflies, key-switching
decompositions, bootstrapping. The configuration is a parameter set. The
values on the wires are polynomials in $\mathbb{Z}_q[X]/(X^N+1)$.

Each field has textbooks.
The composition does not.

\begin{definition}[Reconfigurable homomorphic computing]
The study of compiling \emph{already-synthesized} sequential netlists onto
\emph{tensor engines} so that LUT evaluation, flip-flop clocking, and
(when claimed) torus FHE occupy the same graph object, with explicit
certificates of what the graph does and does not prove.
\end{definition}

HELUT is one laboratory realization of that definition: Yosys JSON in,
\texttt{MPSGraph} out, host posedge clock, LWE/GLWE samples on the FHE
path~\cite{helut-paper}.
The subject is larger than the prototype.
A future chapter may describe a CUDA fabric, an MLIR fabric, or an ASIC LUT
plane. The definition does not name Apple.

\section{Three questions that open the course}

\begin{enumerate}
  \item \textbf{Representation.}
        Can a commodity ML graph API host exact modular arithmetic at
        $q=2^{32}$, or does the API silently round?
  \item \textbf{Clocking.}
        Can sequential logic---enables, sync resets, a RISC-V core---tick
        on that graph without the host becoming a bit-banged FPGA?
  \item \textbf{Honesty.}
        When the same object grows a real bootstrapping key, what
        certificates must ship with every tick, and which sentences remain
        forbidden?
\end{enumerate}

Pillar~I is the constructive answer to (1)--(3).
Pillar~II asks a fourth question the FPGA curriculum never asked:
\emph{what if the INIT table is not binary?}
Pillar~III asks a fifth: \emph{what if the circuit under evaluation is allowed
to mutate because the attacker in the next room is the same compiler?}

\section{Why ``reconfigurable'' is the right word}

An FPGA is reconfigurable because the LUT INIT bits and the interconnect are
late-bound.
HELUT is reconfigurable in the same sense, with three substitutions:

\begin{center}
\begin{tabular}{@{}lll@{}}
\toprule
Classical FPGA & This course & Consequence \\
\midrule
LUT INIT bitstream & Yosys \texttt{\$lut} truth table & Same IR \\
Routing fabric & Tensor edges in one graph & Place-and-route becomes lowering \\
Clock tree & Host \texttt{graph.run} loop & Timing closure becomes tick protocol \\
IOB / BRAM & Placeholders and MTLBuffers & Memory is unified, not dual-port \\
Encrypted? & Optional; certificate-gated & Oracle $\neq$ FHE \\
\bottomrule
\end{tabular}
\end{center}

The last row is the pedagogical load-bearing wall.
Students who have built FPGAs want to say ``the CPU is encrypted'' the moment
PicoRV32 ticks.
The boolean oracle path \emph{does} clock PicoRV32 on ciphertext-\emph{shaped}
vectors (\cid{1}).
The FHE path \emph{does} fetch NOPs on encrypted PicoRV32 at demo $N{=}8$
(\cid{47} CPU, \cid{51} Metal): the PC walks, decrypted $Q$ matches cleartext,
and calibrated hardness is $4.0$ bits---not the production $175.7$ row.
Covering noisy BK at $N{=}1024$ on that core is not claimed.
Chapter~\ref{ch:certs} exists so that confusion is a grading event rather than
a press cycle.

\section{The three pillars as a curriculum, not a brand}

Roadmap language names them the Turing pillar, the Schneier pillar, and a
grand challenge on proving circuits~\cite{helut-trajectory}.
In this course they are simply I, II, III, in the order a student can
verify them:

\begin{itemize}
  \item Without I, II is a float32 netlist toy and III has nowhere to run.
  \item Without II, I is a compiler course with FHE homework.
  \item Without III, I and II only attack the past; they do not design the
        next cipher under fire.
\end{itemize}

Edition~\corpusedition{} teaches I in lecture depth, Pillar~II
Theorem~\ref{thm:tensorlut} (\cid{19}) in lecture depth with empirics in
seminar, Pillar~III in seminar depth, and the grand-challenge / side-channel
items as labeled frontier (Chapter~\ref{ch:frontier}).

\section{Related systems, scoped not ranked}

\begin{center}
\begin{tabular}{@{}lcccc@{}}
\toprule
System & Exact $q{=}2^{32}$ Metal & Netlist clock & BK \texttt{\$lut} FHE & Diff.\ hardware \\
\midrule
Concrete / tfhe-rs & -- & compiler IR & yes & -- \\
OpenFHE / HElib & -- & circuit API & yes & -- \\
HELUT boolean oracle & yes & yes & -- & -- \\
HELUT encrypted path & yes & yes & yes & -- \\
TensorLUT (Pillar~II) & -- & Yosys INIT & -- & yes \\
\bottomrule
\end{tabular}
\end{center}

This is a \emph{scope} matrix, not a throughput bake-off~\cite{helut-paper}.
The encrypted-path ``yes'' is GGSW \texttt{\$lut} FHE as certified
(\cid{4}--\cid{6}, covering \cid{52}--\cid{54}); it is not PicoRV covering
at production $N$.
Concrete is a production TFHE compiler; HELUT is a netlist-clocked tensor
fabric with an FHE path.
Students comparing ``ms/gate'' across those rows without naming encoding,
$N$, and certificate status fail the five-cell test.

\begin{exercise}
Write four sentences, one per row of the matrix, that a journalist could
print without creating a \nid{}.
Then write the one-sentence lie each row most tempts.
\end{exercise}

\growswhen{A fourth pillar chapter, if FHE-gate / ZK-depth work in the
mid-term trajectory~\cite{helut-trajectory} graduates. Until then it is
Chapter~\ref{ch:open}, not a \texttt{\textbackslash part}.}

\chapter{Netlists as the universal IR}
\label{ch:reconf}

\section{From RTL to LUTs}

A combinational Boolean function $f:\{0,1\}^k\to\{0,1\}$ is exactly a truth
table of $2^k$ bits.
FPGA vendors package that table as a $k$-input look-up table.
Yosys~\cite{yosys} emits the same object after \texttt{abc -lut 2} (or 4, or 6)
as a cell of type \texttt{\$lut} with a parameter string INIT.

\begin{definition}[Yosys \texttt{\$lut}]
A combinational cell whose output bit is INIT indexed by the concatenated
input bits. HELUT treats every \texttt{\$lut} as one programmable gate, whether
the backend is a multilinear oracle or a TFHE blind rotate.
\end{definition}

Sequential logic is a \texttt{\$\_DFF*} (or enable / sync-reset variant).
The clock is not a net in the tensor graph: it is a \emph{host protocol}
(Chapter~\ref{ch:pillari}).
That substitution---clock as protocol, not as a routed tree---is the first
reconfigurable-computing heresy of the course, and it is load-bearing.
Tensor engines do not give you a global posedge. They give you
\texttt{graph.run}.

\section{Why we start from synthesized netlists}

High-level FHE compilers start from a language and lower to homomorphic ops.
HELUT starts from a \emph{netlist that already exists}: PicoRV32, a regex DFA,
a decision tree, an Enigma core~\cite{picorv32}.
The pedagogical reason is the same as the systems reason:

\begin{itemize}
  \item The IR is stable. Verilog front-ends change; \texttt{\$lut} does not.
  \item Sequential depth is explicit. Multiplicative depth in an arithmetic
        circuit is a metaphor; DFF count is a fact.
  \item Any HDL, any synthesizer that can emit Yosys JSON, is in scope.
  \item Cryptanalysis targets (historical machines, stream ciphers) already
        \emph{are} netlists.
\end{itemize}

\begin{remark}
Starting from netlists does not make HELUT ``better than Concrete.''
It makes it a different compiler problem: \emph{spatial} Boolean fabric
onto a graph API, not \emph{scalar} program onto PBS.
\end{remark}

\section{The four objects a student must be able to draw}

\begin{enumerate}
  \item \textbf{A LUT} as a box with INIT and $k$ wires.
  \item \textbf{A DFF} as a box whose $Q$ at tick $t+1$ is a mux of $D$, enable,
        and reset, evaluated \emph{between} graph runs.
  \item \textbf{A batch axis $B$} as $B$ independent copies of the same fabric
        in one run---the reconfigurable analog of SIMD, and the reason a
        Bombe search is a tensor problem.
  \item \textbf{A polynomial degree $N$} as the \emph{ciphertext shape}, not
        the circuit size. $N=1$ is clear-shape boolean. $N=1024$ is
        TFHE-shaped. Confusing $N$ with LUT count is the most common
        homework error in week~2.
\end{enumerate}

\section{Worked lowering: a half-adder}

A half-adder is two gates: $\mathrm{sum}=A\oplus B$, $\mathrm{carry}=A\land B$.
Under additive torus encoding, XOR is a vector add (free).
AND is a LUT (or a PBS).

\begin{example}[Half-adder as a graph]
Yosys emits (typically) one \texttt{\$lut} for the AND and uses an add/xor
cell or a second LUT for the sum, depending on techmap.
HELUT's Phase-2 compiler~\cite{helut-paper} lowers this to placeholders for
$A,B$, an add node, a LUT node, and output binds---one persistent
\texttt{MPSGraph}.
There is no FPGA routing step. The ``place-and-route'' \emph{is} the tensor
edge list.
\end{example}

\begin{exercise}
Synthesize a 1-bit full adder with Yosys (\texttt{abc -lut 2}) and count
\texttt{\$lut} versus \texttt{\$\_DFF*} cells.
Then argue, in four sentences, why the encrypted SING of that netlist is a
better first FHE lab than PicoRV32. (PicoRV32 is \cid{1} on the \emph{oracle}
path; the encrypted path's multi-LUT workhorse is the adder, \cid{6}.)
\end{exercise}

\section{Batch is a first-class capacity knob}

On an FPGA, duplicating a circuit costs LUTs.
On a tensor engine, duplicating a circuit is a batch dimension $B$ on every
placeholder.
Memory scales as the broadcast intermediates, not as ``more CLBs.''

\reproduced{C1}{Yosys \texttt{\$lut}/DFF $\to$ Metal/CPU tensor eval on the
cleartext / trivial-oracle path. PicoRV32, Enigma stream, and campaign-scale
batch search all sit on this lowering.}

At $N=1024$, $B=1000$, Enigma~M3's TFHE-shaped working set is gigabytes while
clear-shape $N=1$ stays flat.
Students designing a ``just crank $B$'' attack are designing a memory bomb.
That is a reconfigurable-computing lesson, not an FHE lesson.

\begin{exercise}
Using only the boolean-path scale table in the HELUT README (epoch of this
edition), plot tick time versus $B$ at $N=1$ and $N=1024$.
Identify the regime where graph overhead dominates $N$, and the regime where
$N=1024$ becomes memory-bound.
\end{exercise}

\section{What this chapter refuses to call a bitstream}

A Metal graph executable is late-bound, cached, and device-specific.
It is tempting to call it a bitstream.
We do not, yet.
A bitstream has a documented configuration memory; an \texttt{MPSGraphExecutable}
has an Apple runtime.
The analogy is useful in lecture and false on an exam.
Chapter~\ref{ch:metal} returns to this when the compiler stops unrolling
schoolbook arithmetic into MLIR and starts calling ring kernels.

\chapter{Torus arithmetic and programmable bootstrapping}
\label{ch:torus}

This chapter is a working minimum of TFHE~\cite{chillotti2020tfhe,ducas2015fhew},
not a substitute for a lattice-crypto course.
The goal is that a student can read a HELUT certificate line without bluffing.

\section{The torus and the native word}

\begin{definition}[Discrete torus used here]
HELUT freezes the modulus $q=2^{32}$ and represents torus elements as
\texttt{UInt32}. Addition is native wraparound. This is a \emph{systems}
choice: the machine word \emph{is} the torus.
\end{definition}

That freeze answers a common question: \(q=2^{32}\) is \emph{not} a proof that
lattice FHE requires a 32-bit modulus, and TensorLUT (\cid{19}) is already
Boolean / \([0,1]\). Binary-modulus FHE is a new experiment, not a silent
reread of \cid{4}--\cid{6}. Split: \corpus{directives/q-32-vs-q-2.md}.

The ring of polynomials is
\[
  R_q = \mathbb{Z}_q[X]/(X^N+1),
\]
the negacyclic ring. Multiplication by $X$ is a rotation with a sign flip on
wrap-around---the algebraic fact that makes blind rotation a barrel shifter
on coefficients.

\begin{warning}
Exact $\mathrm{mod}\,2^{32}$ integer tensors are not ``homomorphic encryption.''
They are the arithmetic. Encryption begins when those tensors are LWE or GLWE
\emph{samples} with a secret and noise.
\end{warning}

\section{LWE samples, in one page}

\begin{definition}[LWE sample~\cite{regev2005lwe}]
For secret $s\in\mathbb{Z}_q^n$, a sample is $(a,\langle a,s\rangle+e+m)\in\mathbb{Z}_q^n\times\mathbb{Z}_q$,
with $a$ uniform and $e$ small Gaussian (or discrete) noise.
GLWE is the polynomial analog: mask polynomials plus a body polynomial in $R_q$.
\end{definition}

HELUT's encrypted path uses LWE/GLWE samples and GGSW bootstrap keys.
Decision-LWE is bound to IND-CPA by a machine-checkable certificate
(\cid{5}); the bit estimate at production $(n,\sigma)=(1024,2^{16})$ is
HELUT calibrated $175.7$ versus Sage Cost $180.2$ on that named row
(\cid{23}) and is \emph{not} ``176-bit secure''---print \hid{1} with the number.

\section{Programmable bootstrapping, sketched}

A TFHE Boolean gate is often evaluated by programmable bootstrapping (PBS):
an encrypted phase selects a coefficient of a \emph{test polynomial} via
\emph{blind rotation}, then a sample extract returns an LWE encryption of the
LUT output.

Blind rotation is a product of CMUXes, one per bit of the encrypted phase.
Each CMUX is an external product against a GGSW encrypt of a secret bit.
The accumulator lives in $R_q$ (a GLWE).
The sequential dependence across the $N$ (or $n$) bits is real: you cannot
parallelize away the chain inside one bootstrap.
You \emph{can} refuse to unroll each external product into $O(N)$ host MLIR
nodes. That refusal is Chapter~\ref{ch:metal}.

\section{Negacyclic Toeplitz embedding}

Let $A(X)=\sum_{i=0}^{N-1}a_i X^i$.
Negacyclic convolution against a vector $x$ is $M_A x$, where $M_A$ is the
Toeplitz matrix whose first column is $(a_0,\ldots,a_{N-1})^\top$ and each
subsequent column is the previous column rotated downward with the wrapped
entry \emph{negated}~\cite{lyubashevsky2010ideal}.

\begin{example}[$N=4$ by hand]
Let $A=(a_0,a_1,a_2,a_3)$. Then
\[
M_A
=
\begin{pmatrix}
a_0 & -a_3 & -a_2 & -a_1 \\
a_1 &  a_0 & -a_3 & -a_2 \\
a_2 &  a_1 &  a_0 & -a_3 \\
a_3 &  a_2 &  a_1 &  a_0
\end{pmatrix}.
\]
On \texttt{UInt32}, $-a\equiv 2^{32}-a$. The matrix product \emph{is} the
ring product. HELUT's Phase-1 kernel proved this bit-exact against a CPU
schoolbook oracle, including torus corners $\{0,1,2^{31},2^{32}-1\}$.
\end{example}

\begin{exercise}
Expand $M_A$ for $A=(1,2,3,4)$ over $\mathbb{Z}/2^{32}\mathbb{Z}$ and multiply
by $x=(1,0,0,0)^\top$.
Interpret the output as a polynomial. Then multiply by $x=(0,1,0,0)^\top$ and
explain the sign that appears.
\end{exercise}

\paragraph{Why dense Toeplitz is not the FHE hot path.}
Materializing $N\times N$ at $N=1024$ is $\Theta(N^2)$ words per LUT---a
kernel proof and a memory stress, not a production PBS.
The encrypted hot path moved through tiled schoolbook, fused external product,
GPU-resident tiles, and 3-prime NTT (\cid{16}--\cid{18}).
Students implementing ``the HELUT paper kernel'' as their FHE engine are
implementing 2026-era Phase~1, which the compiler chapter treats as a
\emph{control-plane ancestor}.

\section{Message spacing and \texorpdfstring{$Z_{2N}$}{Z\_2N}}

Blind rotation addresses coefficients at granularity $\delta = q/(2N)$.
Packing several LWE bits into one phase must respect the cyclic group
$Z_{2N}$, not an integer that merely looks like a power.

\reproduced{C6}{Multi-LUT encrypted full\_adder SING PASSes at
$N\in\{256,512,1024\}$; this is what closed \hid{2} on 2026-08-12. The bug was
arity-3 packing overflowing $256\cdot 2N$ headroom; the fix is rotation-native
encrypt plus \texttt{packLWEBits} reduce mod $2N$, with \texttt{rotationPower}
preferring $Z_{2N}$ representatives.}

That closed hedge is the best homework in the chapter: a systems bug that is
also an algebraic bug.

\begin{exercise}
Suppose $N=256$ so $2N=512$. A pack of three bits each scaled by $256$ can
leave the $2N$ range.
Show a numeric counterexample, then write the reduction that restores a
valid rotation power.
\end{exercise}

\section{Noise, in the sense this course needs}

Three layers ship as code, not as vibes (\cid{5}):

\begin{center}
\begin{tabular}{@{}lp{0.62\textwidth}@{}}
\toprule
\texttt{TFHENoiseProof} & Discrete $\infty$-norm inject lemmas. \\
Gaussian $\varepsilon$-certificate & Ingest failure $\le 2^{-64}$ under
independent-noise hypotheses. \\
\texttt{TFHENoisyBKCertificate} & Depth under a BK-noise bound $B_{\mathrm{bk}}$.
Default noiseless Track~A SING still uses $e=0$ BK.
Covering Track~A at $N{=}1024$, torus $\sigma{=}128$, $k{=}7$ is
\cid{52}--\cid{54} (\hid{4} Grade~B).
\texttt{cryptoPublicMS} inject remains \cid{26}/\cid{56}; native $k{=}1$ remains \cid{37}/\cid{55}. \\
\bottomrule
\end{tabular}
\end{center}

\begin{warning}
A noiseless BK is a correct \emph{functional} bootstrap and an incomplete
\emph{depth} story if that is all you ran.
\cid{22} is the covering-gadget measurement at $N\le 128$.
\cid{52}--\cid{54} put covering noisy BK on $N{=}1024$ Metal SING at
$\sigma{=}128$, $k{=}7$.
Students who say ``we have production noisy BK'' from the $N=8$ table, or
who drop the native-$k{=}1$ / \texttt{cryptoPublicMS} remainders, are
failing \hid{4}.
\end{warning}

\section{Parameters you may quote}

From the cookbook~\cite{helut-cookbook}, production-shaped boolean:

\begin{center}
\begin{tabular}{@{}lll@{}}
\toprule
Knob & Value & Asterisk \\
\midrule
$N$ & 1024 & \\
$q$ & $2^{32}$ & native word \\
$\sigma$ & $2^{16}$ & $\ll \delta/2 = 2^{20}$ \\
Target $\varepsilon$ & $\le 2^{-64}$ & Gaussian union on primary wires \\
Classical target & $\ge 128$ bits & HELUT $175.7$ vs Sage $180.2$; \hid{1} \\
$B_{\mathrm{bk}}$ & covering Track~A \cid{52}--\cid{54}; native $k{=}1$ \cid{37}/\cid{55}; \texttt{cryptoPublicMS} \cid{26}/\cid{56} & \hid{4} remainder \\
Inter-LUT refresh & \texttt{publicMS} default & lattice-compatible BK masks \\
\bottomrule
\end{tabular}
\end{center}

Demo $N=8$ will not meet 128-bit hardness. The certificate reports that
honestly. Using demo $N$ for a ``secure encrypted CPU'' slide is a \claimkind{N}.


\part{Pillar I --- Netlist-Clocked Torus FHE}

\chapter{Netlist-clocked torus FHE}
\label{ch:pillari}

Pillar~I is the reason the course exists: a synthesized sequential circuit
becomes one tensor graph, ticked by the host, with an encrypted backend that
is actually LWE.

\section{The HELUT question}

\begin{quote}
\emph{Can Apple Silicon's \texttt{MPSGraph} host exact torus modular arithmetic
and gate-level sequential circuits---and can that same hardware object carry
real bootstrapping keys without lying about it?}
\end{quote}

The first half is a systems existence proof.
The second half is a certificate problem.
HELUT answers both, with the oracle path and the FHE path named as different
objects~\cite{helut-paper,helut-release}.

\section{Pipeline}

\begin{center}
\fbox{\parbox{0.92\linewidth}{\centering\vspace{0.5em}
Verilog $\xrightarrow{\text{Yosys}}$ JSON netlist
$\xrightarrow{\text{HELUTCore}}$ one \texttt{MPSGraph}\\[0.35em]
Host Swift loop $=$ posedge clock; DFF $Q$ ping-pongs across ticks\\[0.15em]
\texttt{\$lut} backend $=$ oracle (multilinear / trivial PBS)
\quad or\quad
encrypted (BK blind rotate)
\vspace{0.5em}}}
\end{center}

All FHE-path arithmetic tensors remain \texttt{MPSDataType.uInt32}.
Float32 belongs to TensorLUT (Pillar~II), not to this chapter.

\section{Compiler lowering}

\texttt{YosysGraphCompiler} walks a \texttt{write\_json} module:

\begin{enumerate}
  \item \textbf{Inputs.} Each input port bit becomes an \texttt{InputNode}
        placeholder of shape $[B,N]$.
  \item \textbf{State.} Each recognized flip-flop allocates a $Q$ placeholder
        \emph{before} LUT lowering so sequential feedback nets resolve.
  \item \textbf{LUTs.} Cells schedule until all drivers exist; combinational
        cycles abort.
  \item \textbf{Next-state.} Enables and sync-resets become exact
        $\mathbb{Z}/2^{32}\mathbb{Z}$ muxes, e.g.\ active-high enable
        \[
          Q_{\mathrm{next}} = E\cdot D + (1-E)\cdot Q.
        \]
  \item \textbf{Outputs.} Driven nets bind to port tensors; Yosys \texttt{"x"}
        becomes constant~0.
\end{enumerate}

Cell coverage includes \texttt{\$\_DFF*}, \texttt{\$\_DFFE*},
\texttt{\$\_SDFF*}, \texttt{\$\_SDFFE*}, \texttt{\$\_SDFFCE*}---the techmap
PicoRV32 actually emits.

\section{Host clock}

Each tick is one \texttt{graph.run}.
With \texttt{retainHistory=false}, two preallocated state buffer sets ping-pong
and output scratch is reused, so a 1\,000-tick stress test does not leak
\texttt{MTLBuffer} retention.
Primary port feeds (including scripted \texttt{resetn}) are host-mutable shared
buffers; LUT matrices or BK packs are uploaded once and cached.

This is the reconfigurable clock tree: not routed silicon, but a protocol
with explicit buffer ownership.

\section{Two backends, one netlist}

\begin{center}
\begin{tabular}{@{}p{0.28\textwidth}p{0.28\textwidth}p{0.32\textwidth}@{}}
\toprule
Backend & Claim status & Use in this course \\
\midrule
Trivial / multilinear / CMUX oracle
  & Shape laboratory, \emph{not} FHE
  & Batch scaling, PicoRV32 boot, Enigma letter clocks. \\
\texttt{encrypted} (LWE + GGSW BK)
  & Pillar~I FHE claim (\cid{4}--\cid{6})
  & SING vs cleartext; certificates on every tick. \\
\bottomrule
\end{tabular}
\end{center}

Inter-LUT refresh defaults to \texttt{publicMS} (lattice-compatible BK masks).
\texttt{.secret} remains for coverage. \texttt{.none} is a debug foot-gun.

\reproduced{C4--C6}{Encrypted LWE/GLWE $+$ GGSW blind-rotate per \texttt{\$lut};
encrypted $\equiv$ clear multi-netlist SING; full\_adder through $N=1024$
(publicMS and secret).}

\section{Why PicoRV32 is the oracle capstone, not the FHE capstone}

PicoRV32 lowers to 4\,785 LUTs and 1\,565 DFFs~\cite{picorv32}.
On the boolean-oracle path it compiles in ${\sim}1.3\,\mathrm{s}$ and ticks
at ${\sim}173\,\mathrm{ms}$ ($N=1024$, $B=1$) in the 2026-08-12 boolean bench
(\cid{1}).
That is a fabric existence proof: a commodity neural/GPU graph can host a
soft CPU's \emph{shape}.

The encrypted path's correctness workhorse is smaller: \texttt{full\_adder},
then tree and regex netlists, under SING.
An encrypted PicoRV32 at production $N$ is a trajectory item, not a
\cid{} row in this epoch.

\begin{warning}
``We booted an encrypted RISC-V core'' is true of the \emph{dark-VM} story
only in the oracle or mock-encoding sense unless a SING log says otherwise.
Avenue~2 in Chapter~\ref{ch:frontier} is the speculative sequel.
Do not merge those sentences on a slide.
\end{warning}

\section{SING as the teaching metric}

SING (stimuli in, golden out) is encrypted $\equiv$ clear on generic vectors,
with certificate lines.
It is the lab equivalent of a scan chain: if SING fails, you do not get to
discuss milliseconds.

\begin{lab}[Encrypted adder SING]
Run the cookbook SING at demo $N$ and at production-shaped $N$
(Appendix~\ref{app:repro}).
Record pass/fail, ms/row, and which certificate types printed.
Then answer: which of those numbers would you be willing to put in an
abstract without an asterisk?
\end{lab}

\section{Threat / fidelity split}

\begin{center}
\begin{tabular}{@{}lp{0.62\textwidth}@{}}
\toprule
In scope for the \emph{datapath} & Exact modular eval; shape $[B,N]$; DFF
routing; Yosys coverage for PicoRV32. \\
In scope for the \emph{FHE claim} & LWE/GLWE $+$ BK; noise / $\varepsilon$ /
hardness / noisy-BK \emph{certificates} as implemented. \\
Out of scope unless measured & Side channels; quantum; estimator-backed
production keys (\hid{1}); noisy BK as a filled product parameter (\hid{4}). \\
\bottomrule
\end{tabular}
\end{center}

The split is intentional.
If the tensor datapath cannot boot a CPU under mock encodings, adding real
TFHE noise only makes the systems problem harder.
If the FHE path cannot SING an adder, a CPU slide is theatre.

\growswhen{Encrypted sequential clocks beyond combinational \texttt{\$lut}
refresh; async resets and latches; an encrypted PicoRV32 SING at production
$N$.}

\chapter{The Metal torus compiler}
\label{ch:metal}

\section{The picture that changed the compiler}

\begin{verbatim}
Yosys $lut  -->  pack LWE  -->  blind rotate (N CMUXes)
                                 |
                                 +-- failure: each external product expands
                                 |    sum_j b_j * X^j a  as O(N) MPSGraph ops
                                 |    --> millions of host MLIR nodes
                                 |    --> encode wall >> run wall
                                 |
                                 +-- design: small graph that *calls* ring kernels
                                     encode cheap; run on Metal; cache
\end{verbatim}

CMUX dependence across bootstrap bits is sequential and real.
The pain was never ``CMUX cannot run on a GPU.''
The pain was \emph{encoding a schoolbook expansion that should never have been
a graph}~\cite{helut-metal}.

\reproduced{fused $N=1024$ DNF}{Schoolbook-in-\texttt{MPSGraph} at $N=1024$
spent $11.6\,\mathrm{h}$ on host encode, never reached the GPU, and died
\texttt{SIGTERM 143}. That log is part of the curriculum: a compiler that
does not finish is a measurement, not an anecdote.}

\section{Doctrine}

SoftBus / ANE remains a graph machine.
We stop mistaking ``unroll all ring math into MLIR'' for ``use the GPU.''
\texttt{MPSGraph} is a \emph{staged IR} and a \emph{scheduler}.
Polynomial multiplication and the external product are first-class GPU
operations.

\section{\texorpdfstring{Phase 1 --- survive $N=1024$}{Phase 1 --- survive N=1024}}

Phase 1 does not change the cryptography. It changes the control plane.

\begin{center}
\small
\begin{tabular}{@{}p{0.22\textwidth}p{0.30\textwidth}p{0.38\textwidth}@{}}
\toprule
Lever & Idea & Teaching point \\
\midrule
Tiled blind rotate & CMUX windows of width $W$; re-ingest ACC
  & Encode cost scales with $W$, not fused $N\times$ schoolbook. \\
Constant CSE & Shared splat banks per tile
  & Graph builders die on \texttt{constantWithScalar} spam. \\
Executable cache & Compile once per $(N,\ell,W,\mathrm{BK})$
  & Trial 1 is cold; ticks $\ge 2$ must hit cache. \\
Telemetry & \texttt{encode} / \texttt{gpu\_run} / \texttt{host\_repack}
  & Silence is not a hang if the log counts tiles. \\
\bottomrule
\end{tabular}
\end{center}

Phase~1 without Phase~2 still leaves $O(W\cdot N)$ encode inside each tile.
Phase~2 without Phase~1 is a kernel drop into an unmeasured fused path.
The laboratory did both, in that order.

\section{Phase 2 --- ring kernels}

\begin{enumerate}
  \item \textbf{Negacyclic poly-mul} as a Metal kernel, bit-identical to CPU
        schoolbook (\cid{18} NTT $\equiv$ schoolbook, 19 tests).
  \item \textbf{Fused external product} (\cid{16}): one
        \texttt{helut\_ggsw\_external\_product} launch per CMUX.
        $N=1024$ microbench $1.043\,\mathrm{s}$/BR.
  \item \textbf{GPU-resident BR tile} (\cid{17}): ACC $+$ BK on device;
        $0.519\,\mathrm{s}$/BR at $N=1024$; $N=64$ at $0.001\,\mathrm{s}$/BR.
  \item \textbf{3-prime NTT persist tile} (\cid{18}): gadget digits in-tile;
        BK in NTT domain; pointwise EP $+$ iNTT $+$ CRT.
        $0.433\,\mathrm{s}$/BR at $N=1024$ (gpu $0.43\,\mathrm{s}$, RSS
        $148\,\mathrm{MiB}$).
\end{enumerate}

Default Metal BR: fused if $N\le 64$, tiled-kernel otherwise (NTT EP inside
tiles). Legacy fused megagraph is \texttt{--metal-br-fused} only, and at
production $N$ it is a museum exhibit.

\section{Numbers a student may quote (epoch \livingepoch)}

\begin{center}
\small
\setlength{\tabcolsep}{4pt}
\begin{tabular}{@{}p{0.30\textwidth}p{0.18\textwidth}p{0.28\textwidth}p{0.14\textwidth}@{}}
\toprule
Path at $N=1024$ & s/BR (micro) & SING / 8 rows & RSS \\
\midrule
Fused schoolbook megagraph & DNF ($11.6\,\mathrm{h}$) & --- & ${\sim}4$\,GiB \\
Poly-mul kernel & 3.645 & --- & --- \\
Fused EP (\cid{16}) & 1.043 & $25.1\,\mathrm{s}$ ($3.14$\,s/row) & --- \\
Persist schoolbook (\cid{17}) & 0.519 & $12.2\,\mathrm{s}$ ($1.52$\,s/row) & 68\,MiB \\
NTT persist (\cid{18}) serial SING & \textbf{0.433} & $15.6\,\mathrm{s}$ ($1.95$\,s/row) & 148\,MiB \\
Wavefront-parallel (\cid{20}) & 0.420 (fused 3-prime) & \textbf{$10.6\,\mathrm{s}$} ($1.33$\,s/row) & --- \\
Crypto $\ell=2$ SING (\cid{21}) & --- & $11.38\,\mathrm{s}$ ($1.42$\,s/row) & --- \\
Covering-b1 $k{=}7$ adder (\cid{52}) & --- & $29.40\,\mathrm{s}/2$ public-ms & --- \\
Covering-b2 $k{=}7$ adder (\cid{57}) & --- & $10.33\,\mathrm{s}/1$ public-ms & --- \\
CPU SING full\_adder ($e{=}0$ BK) & --- & ${\sim}52\,\mathrm{s}$ & --- \\
\bottomrule
\end{tabular}
\end{center}

\reproduced{C20, C21}{Wavefront-parallel independent \texttt{\$lut} BRs:
boolean SING $10.6\,\mathrm{s}/8$ beats persist-schoolbook \cid{17};
noiseless cryptoPublicMS $\ell=2$ SING $11.38\,\mathrm{s}/8$ vs \cid{14}'s $175.6\,\mathrm{s}$.
\cid{21} is $e{=}0$ BK. Noisy \texttt{cryptoPublicMS} is \cid{26}/\cid{56}.}

\hedgebox{H3}{The scheduler, not a faster NTT, is what beat \cid{17} on SING.
Serial in-tile NTT (\cid{18}) still loses boolean SING ($15.6\,\mathrm{s}$).
Sage \hid{1} is independent of these wall-clocks.
Do not quote $0.43\,\mathrm{s}$/BR as a PicoRV covering tick: covering-b1
$\ell{=}32$ is a different gadget tax (\cid{52}).}

\begin{exercise}[Compiler post-mortem]
In 400 words: why did the fused megagraph never reach the GPU?
Your answer must distinguish encode wall from run wall, and must not claim
that ``Metal cannot multiply polynomials.''
\end{exercise}

\begin{exercise}[Reading \hid{3}]
Is \cid{18} a success or a failure?
Defend a one-sentence answer that a claim-audit grader will accept.
(Hint: both, along different bars.)
\end{exercise}

\section{Punchline, for the lecture slide}

Phase~1 makes encoding \emph{finish}.
Phase~2 makes encoding \emph{irrelevant}.
Together they are a graph-native torus compiler: staged IR $+$ ring kernels
$+$ cached executables---instead of a schoolbook megagraph hoping the GPU
will eventually appear.
(SoftBus is Pillar~III's cipher contract, not this Metal lowering.)

\growswhen{NTT EP inside the crypto $\ell=2$ tile (C21's SING log still
shows $W=64$ schoolbook tiles); then \hid{3} is only Sage \hid{1}.}

\chapter{Certificates as first-class artifacts}
\label{ch:certs}

Most compiler courses grade a binary that runs.
This course grades a binary that \emph{refuses} when it cannot certify.

\section{The surface}

\begin{center}
\begin{tabular}{@{}lp{0.58\textwidth}@{}}
\toprule
Certificate & Role \\
\midrule
\texttt{TFHENoiseProof} & Discrete $\infty$-norm inject lemmas. \\
\texttt{TFHEAsymptoticSecurityCertificate} & Gaussian ingest $\varepsilon\le 2^{-64}$. \\
\texttt{TFHELWEHardnessCertificate} & Decision-LWE $\to$ IND-CPA $+$ bit estimate. \\
\texttt{TFHELWECalibration} & Anchor table for the estimator. \\
\texttt{TFHENoisyBKCertificate} & Depth under $B_{\mathrm{bk}}$ / noiseless hypothesis. \\
\texttt{TFHELWEEstimatorProtocol} & External lattice-estimator merge. \\
\bottomrule
\end{tabular}
\end{center}

\texttt{TFHENoise.refuse} if no certificate.
That is not an API flourish. It is the difference between a demo and a claim.

\section{How to read a SING log}

A passing encrypted tick is a bundle:

\begin{itemize}
  \item Functional: ciphertext decrypts to the clear netlist (SING).
  \item Discrete: $\infty$-norm lemmas held on the inject.
  \item Gaussian: ingest $\varepsilon$ under the stated hypotheses.
  \item Hardness: calibrated classical bits, with \hid{1} attached until Sage
        fills the estimator JSON.
  \item Depth: noisy-BK model, with \hid{4} attached while $B_{\mathrm{bk}}=0$.
\end{itemize}

Students who report only milliseconds have not completed the lab.

\section{Calibration is not estimation}

\begin{center}
\small
\begin{tabular}{@{}lrrrr@{}}
\toprule
Label & $n$ & $\sigma$ & Ref bits & HELUT est. \\
\midrule
demo-N8 & 8 & $2^{12}$ & 4 & 4.0 \\
weak-n256 & 256 & $2^{17}$ & 45 & 40.4 \\
mid-n512 & 512 & $2^{16}$ & 95 & 95.5 \\
classic-n630 & 630 & $2^{15}$ & 128 & 129.0 \\
prod-n1024-s16 & 1024 & $2^{16}$ & 170 & 175.7 \\
\bottomrule
\end{tabular}
\end{center}

$q=2^{32}$, binary secret.
The estimator column is PENDING until the cookbook export is filled.

\hedgebox{H1}{Quote ``$\sim 176$ calibrated classical bits at production
$(n,\sigma)$'' only with the clause ``not lattice-estimator attack cost until
Sage fills the pending JSON.''
A midterm bubble that says ``176-bit secure'' without the clause is wrong.}

\section{Five-cell test, applied to certificates}

Certificates are themselves claims.
They need the five cells:

\begin{enumerate}
  \item Proof: the lemma in code (\texttt{TFHENoiseProof}, \ldots).
  \item Table: the calibration table; the cookbook production row.
  \item Metric: $\varepsilon\le 2^{-64}$; bit estimates; $B_{\mathrm{bk}}$.
  \item Examples: demo $N=8$ (honest fail of 128-bit) and production $N=1024$.
  \item Application: refuse-closed encrypt so a netlist tick cannot silently
        drop hardness.
\end{enumerate}

\begin{exercise}
Which cell is missing for noisy BK as a \emph{product} parameter?
Write the experiment that would fill it (\hid{4}: measured $\sigma_{\mathrm{BK}}$
into \texttt{TFHENoisyBKCertificate}).
\end{exercise}

\section{Honesty as an engineering interface}

The non-claims of Chapter~\ref{ch:living} are not ethics theater.
They are API documentation for speech:

\begin{itemize}
  \item Oracle Metal graphs expose the \emph{shape} of encrypted evaluation
        (batch, $N$, clock). They do not expose Decision-LWE.
  \item \texttt{publicMS} is an on-lattice refresh intent. It is not a
        leakage proof (\nid{} on side channels).
  \item A campaign's settings/second is a cleartext Metal number. It does
        not certify encrypted tick rate.
\end{itemize}

\begin{lab}[Certificate literacy]
Run \texttt{.build/release/helut --hardness-table} and classify each printed
row as demo, weak, or production-shaped.
For each row, write one sentence a journalist may print and one sentence
they may not.
\end{lab}

\growswhen{Estimator-backed rows replace the PENDING column; then this
chapter gains a comparison table against Concrete / tfhe-rs parameter sets
as a \emph{scope} table, still not a bake-off unless the bars match.}


\part{Pillar II --- Differentiable Hardware}

\chapter{Continuous hardware, discrete silicon}
\label{ch:pillarii}

The structural theorem of this chapter is corpus-safe (\cid{19}) --- not a
finished lecture.
A plain-English restatement lives in \corpus{directives/theorem-1-plain.md};
a Swift-free check is \texttt{python3 Scripts/tensorlut\_math\_ref.py}, and the
smallest worked example of the LUT view is
\texttt{python3 Scripts/toy\_cipher\_demo.py} (\corpus{INTRO.md}).
The separable melt--freeze--snap certificate is corpus-safe (\cid{44}).
Shatter / hold grades and campaign empirics remain seminar-depth and are
\emph{not} the theorem~\cite{helut-tensorlut-thm}.

\growswhen{A completeness theorem for \emph{multi-LUT topological} melt:
that melt recovers cryptographic structure combinatoric sieves miss, for a
stated class of netlists beyond a fully observed 1-LUT.
Theorem~1$''$ (\cid{44}) is the separable interpolant, not that class.
Stream-cipher melts remain on the trajectory.}

\section{The problem boolean search cannot see}

A combinatoric sieve on a cipher netlist proposes discrete genotypes:
INIT bits, plugboard pairs, rotor orders.
On a short ciphertext the search space is both enormous and full of
linguistic hallucinations (the P1030680 campaign is the case study, not
the claim).
Boolean search \emph{rejects}.
It does not \emph{propose} a nearby reciprocal structure that is almost a
stecker.

Pillar~II melts the hardware~\cite{helut-paper,helut-tensorlut-thm}.

\section{Setup}

Let $L$ be the number of LUT6 cells and $w\in[0,1]^{64L}$ the concatenated
INIT genome.
Let $y(w)$ be the \emph{multilinear extension} of those INITs on $[0,1]$
(exact on $\{0,1\}$ inputs).
Let $t$ be a target bit vector of matching width.
Soft crypto fitness and discreteness penalty:
\begin{align}
  F_{\mathrm{crypto}}(w)
  &= -\lVert y(w)-t\rVert_2^2, \\
  \pi(w)
  &= \sum_i w_i(1-w_i), \\
  F(w)
  &= F_{\mathrm{crypto}}(w)-\lambda\,\pi(w),
  \qquad \lambda\ge 0.
\end{align}

\begin{remark}[Implementation schedule]
HELUT's GA uses a squeeze $\lambda=\lambda_{\max}(p')^2$.
Theorem~1 needs only $\lambda\ge 0$.
\end{remark}

Emitter: $E(w)_i=\mathbf{1}[w_i\ge\tfrac12]$.
Freeze mask $M\subseteq\{1,\dots,64L\}$ removes frozen coordinates from
$\pi$.
Stecker genotypes are \emph{partial involutions} (disjoint pairs) on
$\{0,\dots,25\}$.

Every Yosys \texttt{\$lut} of width $1\ldots 6$ pads to a 64-wide
\texttt{Float32} INIT.
The forward pass is branchless and GPU-shaped.
It is not torus FHE.

\section{Theorem 1 (structural)}
\label{sec:thm-tensorlut}

\reproduced{C19}{TensorLUT continuous$\to$discrete Theorem~1:
six machine-checked lemmas
(\texttt{TensorLUTFormal.certificate()}).
Not a cryptanalytic break.
Statement: \corpus{directives/tensorlut-theorem.md}.
Reproduce:
\texttt{swift test -c release --filter testTensorLUTFormalCertificate}.}

\begin{theorem}[TensorLUT continuous$\to$discrete; \corpusthm{1}]
\label{thm:tensorlut}
Assume the hypotheses on the certificate: the LUT is multilinear; the GA
mutates only unfrozen $w_i\in[0,1]$; the involution sandwich freezes core
INITs.
Then:
\begin{enumerate}
  \item \textbf{Discreteness.}
        $\pi(w)\ge 0$, with equality iff $w\in\{0,1\}^{64L}$.
  \item \textbf{Crypto MSE.}
        $F_{\mathrm{crypto}}(w)\le 0$, with equality iff $y(w)=t$.
  \item \textbf{Combined objective.}
        $F(w)\le F_{\mathrm{crypto}}(w)$; if $\pi(w)>0$, increasing
        $\lambda$ strictly decreases $F$.
  \item \textbf{Emitter.}
        $E$ is idempotent on $\{0,1\}^{64L}$ and agrees with threshold
        $\tfrac12$.
  \item \textbf{Involution sandwich.}
        Overlapping pairs are rejected; applying a valid pair-set twice
        is the identity on the alphabet.
  \item \textbf{Freeze.}
        Coordinates in $M$ do not contribute to $\pi$.
\end{enumerate}
\end{theorem}

\paragraph{Proof (machine-checked).}
Each clause is a \texttt{TensorLUTFormal.check*} in
\texttt{TensorLUTFormal.swift}, aggregated by
\texttt{TensorLUTFormal.certificate()}.

\reproduced{C25}{Theorem~1 corollary: emitter--discrete agreement
($\pi=0\Rightarrow E(w)$ recovers INIT bits) and involution under freeze
(\texttt{mutatedPreserving} retains frozen pairs).
\texttt{TensorLUTFormal.corollaryCertificate()}.
Reproduce:
\texttt{swift test -c release --filter testTensorLUTFormalCorollaryCertificate}.}

\begin{theorem}[TensorLUT emitter / freeze corollary; \corpusthm{1} corollary]
\label{thm:tensorlut-corollary}
Under Theorem~\ref{thm:tensorlut} hypotheses plus emit-via-$E$ and
\texttt{mutatedPreserving}:
if $\pi(w)=0$ then $E(w)$ recovers the binary INIT;
every freeze-preserving genotype remains a partial involution containing
the frozen pairs.
\end{theorem}

Still not melt completeness for arbitrary netlists.
The named lemmas of Theorem~\ref{thm:tensorlut} are
\texttt{discretenessPenalty},
\texttt{cryptoFitnessMSE},
\texttt{combinedObjective},
\texttt{emitterThreshold},
\texttt{involutionSandwich},
\texttt{freezeMask}.
Source path: \corpus{Sources/HELUTCore/TensorLUTFormal.swift}.

\begin{warning}
Theorem~1 does \emph{not} prove recovery of arbitrary keys; a U-534 /
P1030680 plaintext; or that melt is complete for all netlists.
Shatter / hold grades remain empirical (\cid{8}, \cid{9}, \hid{6}).
\end{warning}

\reproduced{C44}{Theorem~1$''$ melt--freeze--snap on a coordinate-separable
Boolean interpolant: unique maximizer $w=t$, snap basin
$\lvert w_i-t_i\rvert<1/2$, wrong freeze blocks $F=0$.
\texttt{TensorLUTFormal.meltFreezeSnapCertificate()}.
Reproduce:
\texttt{swift test -c release --filter testTensorLUTMeltFreezeSnapCertificate}.
Not multi-LUT topological melt.}

\reproduced{C48}{Two-LUT cascade $y=(a\land b)\oplus c$: melt from $0.5$,
snap, emit two LUT6 Verilog cells. Snapped INIT matches all 8 Boolean
corners (a dual interpolant such as NAND$+$XNOR is allowed).
\texttt{testTwoLUTCascadeMeltFreezeSnapEmit}. Not Yosys re-synth.}

\begin{theorem}[TensorLUT melt--freeze--snap, separable interpolant; \corpusthm{1$''$}]
\label{thm:tensorlut-snap}
Under a fully observed 1-LUT (each used INIT address an independent Boolean
target $t$), $F(w)=-\lVert w-t\rVert_2^2-\lambda\pi(w)$ has unique maximizer
$w=t$; the emitter recovers $t$ on the open cube
$\lvert w_i-t_i\rvert<1/2$; a freeze away from $t$ makes $F=0$ unreachable.
\end{theorem}

\section{Five-cell test for Theorem 1}

\begin{center}
\small
\begin{tabular}{@{}lp{0.68\textwidth}@{}}
\toprule
Cell & Artifact \\
\midrule
Proof & Theorem~\ref{thm:tensorlut} $+$ \texttt{TensorLUTFormalCertificate} \\
Table & Lemmas vs \texttt{holds} on the certificate; campaign grades in
        \corpus{BREAK\_P1030680.md} §21 \\
Metric & $F_{\mathrm{crypto}}=0$ on unmutated baseline; $\pi=0$ on binary INIT \\
Examples & M4 baseline emit; freeze-core involution; blind 3-pair PASS \\
Application & Stecker search that cannot propose a non-reciprocal map \\
\bottomrule
\end{tabular}
\end{center}

\begin{exercise}[Read a certificate]
Run the reproduce filter (Appendix~\ref{app:repro}) or, on the theory track,
read \texttt{TensorLUTFormal.certificate()} and list the six
\texttt{holds} bits.
For each lemma, write one sentence a journalist may print and one they
may not.
\end{exercise}

\section{Shatter versus hold (empirical)}

A full cold-start on a large sequential cipher discovers continuous
shortcuts that \emph{shatter} when $\lambda$ rises: $F_{\mathrm{crypto}}$
looks perfect in float and dies on the Boolean cube.
Targeted melting freezes known-good LUT blocks (\texttt{freezeMask}) and
evolves only the unknown cone---clause~6 of Theorem~\ref{thm:tensorlut}
is why frozen coordinates do not pay $\pi$.
Early resignation reboots doomed lineages rather than squeezing a
reciprocal shortcut to the end.

\begin{definition}[Grade]
Under a published $\lambda$ schedule, a lineage \textbf{holds} if the
emitted binary netlist preserves $F_{\mathrm{crypto}}=0$ (or a stated
bar). It \textbf{shatters} if discreteness kills the continuous solution.
Both grades are science. Shatter is not a failed demo; it is evidence
about shortcuts. It is \emph{not} a clause of Theorem~\ref{thm:tensorlut}.
\end{definition}

\reproduced{C8, C9}{TensorLUT M4 baseline $F_{\mathrm{crypto}}=0$ plus Verilog
emit; Stecker involution sandwich with blind 3-pair PASS.
Method claims, not a U-534 / P1030680 plaintext.}

\section{Involution sandwich}

Stecker genotypes are partial involutions (disjoint pairs) by construction
(Theorem~\ref{thm:tensorlut}, clause~5).
The GA cannot propose a non-reciprocal map.
The 16 cone LUTs around an identity plugboard cannot invent letter swaps by
INIT melt; the sandwich injects $S(\mathrm{CT})$ into a frozen core and
scores soft PT against bits of $S(P)$.
Reciprocity is structural.

\begin{warning}
Short cribs leave unused pairs unconstrained and can admit false $F=0$
steckers. Grade active-map agreement on $\mathrm{CT}\cup\mathrm{PT}$, use
parsimony, grow the pair budget on plateau, and allow soft freeze $+$ thaw.
TensorLUT grades are not a campaign decrypt (\claimkind{N}, \hid{6}).
\end{warning}

\section{Adversarial synthesis as a compiler loop}

Forward: silicon $\to$ uniform float tensors.
Reverse: cooled tensors $\to$ INIT hex $\to$ Verilog LUT instantiations.
The loop is a generative compiler in \emph{this} laboratory: cooled tensors
emit LUT6 Verilog, not a bitstream onto a physical FPGA.
Theorem~\ref{thm:tensorlut} is the invariant the loop is forced to obey.
It is not a guarantee that the loop finds a key.

\begin{exercise}
Why must $\pi(w)=\sum w_i(1-w_i)$ vanish on binary INITs, and why is a
fitness that ignores $\pi$ allowed to ``solve'' a cipher that no FPGA can
implement?
Give a two-LUT numerical cartoon.
Which clause of Theorem~\ref{thm:tensorlut} did you just unpack?
\end{exercise}

\begin{exercise}[Seminar]
Theorem~\ref{thm:tensorlut} is structural, not complete.
Write the statement of a completeness theorem you wish were true (melt
recovers $X$ on netlist class $Y$ under schedule $Z$), and the
counterexample shape you would try first.
Do not upgrade shatter/hold grades into that theorem.
\end{exercise}

\nonclaimbox{Pillar~II does not break P1030680, does not replace Welchman
search, and does not make TensorLUT a PBS. Parallel research; campaign
fitness stays cleartext.
\cid{19} is a theorem about the objective and the involution contract,
not about plaintext.}


\part{Pillar III --- Adversarial Polymorphic Ciphers}

\chapter{Ciphers that grow under fire}
\label{ch:pillariii}

\growswhen{A polymorphic Red/Blue standard beyond the bounded E256 fixture-v4
receipt (roadmap: Schneier pillar), and IND-CPA / Red-pressure security results
that go past finite structural checks. This chapter is a scaffold and collecting
surface for existing receipts---not lecture-ready material and not a finished
industry standard.}

\section{Evidence boundary for a design laboratory}

Historical Enigma traits motivate the E256 design laboratory, but the checked
claims here do not prove those traits universally fatal or prove their removal
for every generated machine. This collecting surface records three existing
rows: bounded fixture-v4 implementation parity (C10), a bounded structural
certificate / protocol (C24), and a historical encrypted teaching cone (C39).

The same-state reciprocal shape is a design target. The evidence below is
finite: only checks explicitly called exhaustive are exhaustive, and key/state
coverage remains bounded.

\paragraph{Open E256-H/X0 hardware and policy measurements.}
The E256-061/062/063 hardware and X0 checks remain
\texttt{OPEN\_PROGRESS} and have no C/H/N row. Their complete dependency-ordered
reproduction commands, receipt digests, and non-claims are in
\corpus{REPRODUCE.md}. The three latest gates can be re-checked directly with:
\begin{itemize}[nosep]
  \item \texttt{make e256-hw-h2-check};
  \item \texttt{make e256-hw-h4-sequence-check};
  \item \texttt{make e256-x0-h4-evolution-check}.
\end{itemize}
The H4 census validly found that 737,280 of 1,474,560 preregistered mixed
three-round windows fail the exact all-to-all byte-dependency criterion. Its
first frozen-order witness uses list indices $[1,2,3]$---tuples
$(0,1,3,4)$, $(0,1,4,3)$, and $(0,1,4,7)$---where output lane 0 lacks input
lanes $[6,11,24,29]$. The sequence-specific H3 collapse passed 442,368 exact
comparisons; integrated RTL/model and per-input cycle evidence passed 55,296
comparisons each; and all 19 controls were detected.

X0 then checked the exact criterion $S(y)+S(z)=\mathbb{Z}_8$ for all 147,456
ordered catalog pairs using independent literal and sumset implementations.
Exactly 73,728 directed transitions are safe and every node has in/out degree
192. The full graph has two 192-node strongly connected components; the
16-support quotient has two 8-node components and no Hamiltonian cycle. The
frozen fallback certificate alternates catalog indices 0 and 3 for twelve
selectors and verifies all 11 adjacent transitions. All 23 X0 controls were
detected. This proves only the represented contiguous three-round structural
dependency windows; the two components prevent one safe walk from covering all
16 supports.

The generic Yosys/ABC sequence-minus-repeated observations were:
\begin{itemize}
  \item FF storage: $-407$ LUT6, $+132$ state bits, $-275$ primitive
        cells, and $+1$ combinational level;
  \item ring storage: $-57$ LUT6, $+132$ state bits, $+75$ primitive
        cells, and $+1$ combinational level.
\end{itemize}
There was no threshold or winner. Negative LUT deltas from separate generic
mappings are not physical selector savings. The H4 counterexample parks
unrestricted arithmetic switching; X0 supplies a bounded safe-transition
grammar and path, not a production evolution mechanism. Neither result inherits
the fixed-tuple four-round/25-active-S-box statement. These are generic
mapped-RTL measurements, bounded falsification, and host-side structural
evidence only---not a theorem or course-ready laboratory result. They select no
architecture, storage organization, round count, schedule, XOF, H4 policy,
suite, protocol, production RTL, cryptanalytic result, or security value.

\section{Fixture-v4 evidence capsule (scaffold)}

\begin{itemize}
  \item The fixed live profile is schema-4 \texttt{E256-v2/gen0}, fixture-v4.
  \item Its conjugated-XOR center is
        $A_i^{-1}(A_i(x)\oplus k_i)$.
  \item The host derives and transports
        \texttt{(payload, centerMask, absoluteByteCounter)}; RTL validates
        \texttt{absoluteByteCounter} and does not implement HMAC.
  \item The formal C24 test is \textbf{1/1 PASS}; the full E256 C10 suite is
        \textbf{49/49 PASS}. The receipt is
        \texttt{logs/e256-v2-gen0-fixture-v4-validation.json}.
\end{itemize}

\reproduced{C10}{Enigma256 fixture-v4 bounded implementation / KAT-parity,
reciprocity, and conjugated-XOR evidence for the fixed schema-4
\texttt{E256-v2/gen0} native profile. Full E256 suite: \textbf{49/49 PASS}.
Receipt: \texttt{logs/e256-v2-gen0-fixture-v4-validation.json}.}

\reproduced{C24}{Enigma256 SoftBus Theorem~2 is a bounded fixture-v4
structural certificate / protocol: deterministic four-state
scramble-bijection sweep; exhaustive byte reciprocity for one frozen fixture
state; deterministic 128-byte stream round-trip; fixture plugboard plus
selected XOR-center involution/fixed-point checks; and exact schema-4
\texttt{E256-v2/gen0} native-profile integrity. Formal test:
\textbf{1/1 PASS}. Statement: \corpus{directives/enigma256-theorem.md}.}

\reproduced{C39}{One frozen historical one-byte / one-round E256 teaching cone
with algebraic sboxes, legacy UKW, offsets $0$, and identity plug under GGSW
Metal SING at $N{=}1024$. It predates the live fixture-v4 profile and does not
exercise transported mask/counter tuples, the conjugated-XOR center, live BRAM
day-key tables, NLFF stepping, or the full core.}

\section{Theorem 2 (bounded structural certificate / protocol)}
\label{sec:thm-enigma256}

\begin{theorem}[E256 fixture-v4 bounded structural certificate / protocol;
\corpusthm{2}]
\label{thm:enigma256}
For the fixed schema-4 \texttt{E256-v2/gen0} native-profile fixture-v4, the C24
bounded structural certificate / protocol checks exactly: (i) a deterministic
four-state scramble-bijection sweep; (ii) exhaustive byte reciprocity for one
frozen fixture state; (iii) a deterministic 128-byte stream round-trip;
(iv) the fixture plugboard plus selected XOR-center involution/fixed-point
checks, where the center is $A_i^{-1}(A_i(x)\oplus k_i)$; and (v) exact
schema-4 \texttt{E256-v2/gen0} native-profile integrity. The host derives and
transports \texttt{(payload, centerMask, absoluteByteCounter)}. RTL validates
\texttt{absoluteByteCounter} and does not implement HMAC.
\end{theorem}

Machine-checked result: \texttt{Enigma256Formal.certificate()} /
\texttt{testEnigma256FormalCertificate}, \textbf{1/1 PASS}; full E256 suite,
\textbf{49/49 PASS}. Receipt:
\texttt{logs/e256-v2-gen0-fixture-v4-validation.json}.
Reproduce through \corpus{REPRODUCE.md}.

\begin{warning}
Theorem~2 records a \emph{bounded structural certificate / protocol}, not a
universal proof. Finite checks are exhaustive only where stated, and state/key
coverage is bounded. It is not IND-CPA, an HMAC proof, external cryptanalysis,
human acceptance, or a finished polymorphic-cipher standard.
\end{warning}

\section{Red-pressure questions kept separate}

TensorLUT, known-plaintext, and \texttt{ent} experiments remain internal Red
questions on this collecting surface. C24 does not show that those experiments
cannot force a generation roll, and the fixture-v4 receipt is not external
cryptanalysis or human acceptance. No coupled-NLFF hardening result is part of
the bounded certificate.

\section{Draft lab prompts (scaffold; not lecture-ready)}

\begin{exercise}[Blue studio prompt]
Propose one additional E256 invariant as an open design question (for example,
an NLFF period lower bound or BRAM burst-width condition). Specify a test that
could refuse a generation, but do not promote the proposal to a claim without
a checked receipt.
\end{exercise}

\begin{exercise}[Historical-leak checklist prompt]
Compare historical Enigma traits with the corresponding E256 design choices.
Keep every unsupported security consequence labeled as a question, not a
claim.
\end{exercise}

\section{Relation to Pillars I and II}

This collecting surface links the bounded C24 protocol to its fixture receipt
(Theorem~\ref{thm:enigma256}). \cid{39} is a historical frozen one-round
legacy-UKW scramble under GGSW Metal SING. It predates live fixture-v4 and does
not exercise transported \texttt{centerMask}/\texttt{absoluteByteCounter}, the
conjugated-XOR center, live BRAM, NLFF, or a full core. Calling that lab
``production encrypted E256'' is a \claimkind{N}.

\nonclaimbox{E256 is not a NIST candidate in this scaffold, not a replacement
for TLS, and not an implication that polymorphic rotors beat lattices. The
fixture-v4 receipt is bounded implementation / protocol evidence, not IND-CPA,
an HMAC proof, external cryptanalysis, or human acceptance.}


\part{Laboratory, Applications, and the Frontier}

\chapter{Application gallery}
\label{ch:apps}

Applications exist to keep pillars honest.
A compiler with no netlist is a speech; a melt with no cipher is a float
toy; a polymorphic cipher with no Red is a blog post.

\section{Tiered datapath (Pillar~I shape laboratory)}

\begin{enumerate}
  \item \textbf{Decision tree.}
        A 4-bit non-linear boundary, 7 LUTs. Exact threshold logic on a
        batch of records. Teaches: non-linear Boolean $\neq$ neural net,
        and $B$ is a capacity knob.
  \item \textbf{Regex / DFA search.}
        3-character matcher, 23 LUTs. Large $B$ turns unified memory into
        the story. Teaches: tensor expansion can be a 40\,GiB intermediate
        before it is a clever algorithm.
  \item \textbf{PicoRV32.}
        ${\sim}4.8\mathrm{k}$ LUTs, ${\sim}1.5\mathrm{k}$ DFFs. Scripted
        \texttt{resetn} boot. Teaches: sequential coverage (enables,
        sync-reset, \texttt{"x"} bits) is the compiler, not the kernel.
        Oracle path: \cid{1}. Encrypted CPU SING: not a \cid{} in this epoch.
\end{enumerate}

\section{Encrypted correctness netlists}

\texttt{full\_adder}, tree, and regex under \texttt{--bench-encrypted --sing}
are the FHE teaching set (\cid{4}--\cid{6}).
The adder is the first lab; the others prove the compiler is not
adder-specialized.

\section{Campaign as a case study (not a decrypt)}

Welchman-path break of known P1030684 (\cid{2}), $\le 10$-plug SAT kill
chain (\cid{3}), Potsdam/Plaice $\neq$ P1030680 (\cid{11}), windowed
discriminator versus whole-message turnover (\cid{12}): these belong in a
cryptanalysis seminar attached to Pillar~II methods.

\nonclaimbox{This book does not claim a decrypt of P1030680.
Campaign fitness is cleartext Metal batch, not HELUT encrypted tick rate.}

\section{E256 (Pillar~III)}

The live cipher under Red pressure (Chapter~\ref{ch:pillariii}).
Use it as a design studio, not as a messaging homework.

\begin{exercise}
Pick one application tier and write a five-cell card for it.
If you cannot name two runnable instances, it is not ready for the gallery
lecture.
\end{exercise}

\growswhen{Figures for the site gallery; an encrypted tree/regex Metal SING
table at production $N$; any new application that earns a \cid{} row.}

\chapter{Laboratories}
\label{ch:lab}

Systems-track labs assume macOS 14$+$, Apple Silicon, Swift 6.3, and a
release \texttt{helut} at repo root:
\texttt{swift build -c release --product helut}.
Theory-track substitutes are stated per lab.

\section{Lab 0 --- Claim literacy (everyone)}

Read \corpus{directives/claim-sheet.md} at the printed epoch.
Classify five public sentences (instructor-provided) as \cid{}, \hid{},
\nid{}, or frontier.
No machine required.

\section{Lab 1 --- Negacyclic oracle}

Implement $N=8$ schoolbook mul in $\mathbb{Z}/2^{32}\mathbb{Z}[X]/(X^N+1)$
and check a dense Toeplitz expansion against it on torus corners.
Theory track: $N=4$ by hand (Chapter~\ref{ch:torus} exercises).

\section{Lab 2 --- Netlist clock}

Compile \texttt{counter} or \texttt{full\_adder} Yosys JSON with
\texttt{--compile-only}.
Tick with \texttt{--bench}. Identify $B$, $N$, LUT count, DFF count.
Explain in writing why $N$ is not the circuit size.

\section{Lab 3 --- Boolean oracle versus FHE path}

\begin{lstlisting}
.build/release/helut --bench netlist.json --degree 64 \
  --lut-backend pbs --encoding phase --bench-equiv
.build/release/helut --bench netlist.json --degree 8 \
  --bench-encrypted --cpu-only --sing --vectors 8
\end{lstlisting}

Deliverable: a table with columns \{backend, $N$, SING, certificates
printed, may we say ``encrypted'' on a slide?\}.

\section{Lab 4 --- Production-shaped SING}

\reproduced{C6, H2 closed}{full\_adder encrypted SING through $N=1024$.}

Follow Appendix~\ref{app:repro}. Record ms/row at $N=8$ and $N=1024$.
Write the \hid{1} and \hid{4} asterisks next to any hardness or depth
sentence.

Theory track: parse a published SING log and reconstruct the table
without running Metal.

\section{Lab 5 --- Metal compiler post-mortem}

Compare fused-megagraph DNF, \cid{16}, \cid{17}, \cid{18} using
Chapter~\ref{ch:metal}'s table.
Optional systems: \texttt{make test-metal-p1} and one
\texttt{--bench-encrypted-micro --degree 64} run.
Deliverable: the \hid{3} paragraph in your own words.

\section{Lab 6 --- TensorLUT Theorem 1 (\cid{19})}

\begin{lstlisting}
swift test -c release --filter testTensorLUTFormalCertificate
\end{lstlisting}

Deliverable: name the six lemmas, the certificate hypotheses, and one
sentence that Theorem~\ref{thm:tensorlut} does \emph{not} prove.
Theory track: read \corpus{Sources/HELUTCore/TensorLUTFormal.swift} and
\corpus{directives/tensorlut-theorem.md}; reconstruct the five-cell card
without running Swift.

\section{Lab 7 --- Five-cell audit (everyone)}

The signature assignment (Chapter~\ref{ch:instructor}).
One abstract sentence from \corpus{paper/helut.tex}; five cells; a
pass/fail on whether the sentence is lecture-safe at this epoch.

\section{Lab 8 --- Frontier proposal (everyone)}

Chapter~\ref{ch:frontier}: pick one avenue. Propose the first experiment
that would mint a \cid{} row. You may not claim the avenue works.

\section{Optional project slots}

TensorLUT involution sandwich replay (\cid{9} logs), E256 invariant
studio, or a new Yosys netlist (student RTL $\to$ SING).
Instructor approval required; campaign decrypt projects are out of scope.

\growswhen{Autograded SING harnesses for course VMs; a Dockerized theory
track; non-Apple ring-oracle labs.}

\chapter{Open problems (trajectory)}
\label{ch:open}

This chapter is the discovery path after disclosure~\cite{helut-trajectory}.
None of these items are lecture facts until they have receipts.

\section{Near term}

\begin{center}
\small
\setlength{\tabcolsep}{4pt}
\begin{tabular}{@{}p{0.24\textwidth}p{0.34\textwidth}p{0.32\textwidth}@{}}
\toprule
Track & Status at epoch \livingepoch{} & Next experiment \\
\midrule
\hid{2} full\_adder $N\ge 256$
  & \textbf{Closed} ($Z_{2N}$ pack / \texttt{rotationPower})
  & Keep as a worked bug in Chapter~\ref{ch:torus}. \\
\hid{1} Sage lattice-estimator
  & \cid{23} filled; production $|\Delta|=4.5$; core-SVP vs Cost \texttt{rop} divergences
  & Optional retune / quote estimator-only on $\Delta>16$ rows. \\
\hid{3} Metal BR at large $N$
  & \cid{20} boolean $10.6\,\mathrm{s}$; \cid{21} crypto $\ell=2$ $11.38\,\mathrm{s}$
  & NTT inside crypto $\ell=2$ at $N{=}1024$ (incomplete public-MS gadget). \\
\hid{4} Noisy BK
  & \cid{52}--\cid{54} covering-b1; \cid{57} covering-b2 $\varepsilon$ + regex SING @ $N{=}1024$
  & Native $k{=}1$ (\cid{37}/\cid{55}); \texttt{cryptoPublicMS} (\cid{26}/\cid{56}); PicoRV covering (\cid{60} Q FAIL). \\
Encrypted sequential
  & \cid{53}/\cid{54} covering-b1 counter + toy ISA @ $N{=}1024$ $\sigma{=}128$ $k{=}7$; \cid{51} PicoRV Metal still demo $N{=}8$
  & Metal PicoRV covering @ $N{=}1024$ (LUT-tax). \\
Campaign catalog
  & Middle ring $\neq A$ untested; catalog parked @417
  & Resume \texttt{--bombe-from 418}. \\
Garble / quarantine
  & Soft-band grades
  & Sister-message lessons; not a decrypt claim. \\
Audience R1--R6
  & Shipped, not \claimkind{C} rows: Linux checks, weak-vs-affine toy cipher pair, Theorem~1 in English, the $q$-split, four-page note
  & Open: Lean/Coq version of the six clauses; toy cipher through Yosys; $q=2$ FHE. \\
Vertex maximizers
  & Literature-resolved, not a \claimkind{C} row
  & Raising $\lambda$ is the classical concave exact penalty for 0--1 programming
    (Raghavachari 1969; Giannessi--Niccolucci 1976). Sufficient:
    $\lambda\ge\frac12\sup\lambda_{\max}(\nabla^2 f)$, measured $2+\sqrt3$ on two LUTs.
    Open: scaling with depth / fan-out. \corpus{note/lut-relaxation.tex} §4--§5. \\
\bottomrule
\end{tabular}
\end{center}

\section{Mid-term pillars}

\begin{enumerate}
  \item \textbf{Pillar I science:} shallower nets, NTT/persist graphs,
        estimator-backed params (this book's Parts~III).
  \item \textbf{Pillar II science:} Theorem~1 landed (\cid{19}); corollary
        emitter / freeze (\cid{25}); melt--freeze--snap on a separable
        interpolant (\cid{44}); 2-LUT cascade emit (\cid{48}); remaining:
        multi-LUT topological melt, stream-cipher melts.
  \item \textbf{Pillar III science:} SoftBus Theorem~2 landed (\cid{24},
        Chapter~\ref{ch:pillariii}); remaining: polymorphic Red/Blue standard
        beyond E256 SoftBus.
  \item \textbf{FHE gate / ZK depth:} TensorLUT aimed at multiplicative
        depth (queued; not a part of this edition).
  \item \textbf{Side-channel (parked):} bgpucap-style power on live Metal
        graphs only with controlled fixtures --- Chapter~\ref{ch:frontier},
        avenue 5.
\end{enumerate}

\begin{exercise}[Pick a next experiment]
Choose one open hedge. Write a one-page lab proposal: command line, pass
bar, which box in this book would change, and which \claimkind{N} you refuse to
imply if the experiment fails.
\end{exercise}

\growswhen{Each closed row in this table deletes a hedge box elsewhere
and adds a reproduced box. That is the only way this chapter shrinks.}

\chapter{Frontier (not claims)}
\label{ch:frontier}

The text in this chapter is adapted from the project's speculative
avenues~\cite{helut-avenues}.
It is included so the public story stays honest: \textbf{reproducible core
now; weird frontier labeled as frontier.}

Instructors: last two seminars, or a 10\% essay.
Never a true/false bank.

\section{Encrypted LLM guardrails}

\frontierbox{Evolve a tiny neural or logic-gate firewall in the encrypted
domain that sits between an untrusted prompt and a local model.
The firewall circuit would evaluate in the dark on token tensors and mutate
its Boolean topology from adversarial inputs it blocked---without decrypting
prompt, weights, or mutation rules in memory.}

First graduating experiment (suggested): a 4-LUT encrypted classifier SING
on synthetic token bits, with a certificate, \emph{before} any talk of
genetic mutation or Ollama.

\section{Self-modifying encrypted binaries}

\frontierbox{A PicoRV32 that already ticks on an \texttt{MPSGraph} under
oracle encodings (\cid{1}) is not a dark VM.
Encrypted PicoRV is demo $N{=}8$ at $4.0$ bits (\cid{45}--\cid{51}), not
production $N$.
The avenue is a binary that rewrites its own instructions while remaining
FHE-encrypted, so the host sees tensor multiplies and cannot inspect the
guest.}

First graduating experiment: encrypted instruction-fetch SING on a toy
ISA smaller than PicoRV32, with a written ethics note. This is a privacy
architecture problem, not a malware homework.

\section{Honey-token ledgers}

\frontierbox{Batched encrypted pattern search composed with a
zero-knowledge ledger, so a sync evaluates trapdoor circuits without the
storage node learning which file or which trigger fired.}

First graduating experiment: regex netlist SING at modest $B$ with an
explicit non-claim that the node learns nothing---then measure what the
\emph{host process} still sees (lengths, timings). Honesty before mesh
networks.

\section{Genetic hardware synthesis}

\frontierbox{\texttt{fhe-evolve} mutates truth tables and connectivity;
Yosys recompiles; \texttt{MPSGraph} re-lowers; the running system processes
a stream through hardware it just invented.}

First graduating experiment: one offline generation of a 2-LUT netlist that
SING-matches a fixed spec, wall-clock of synth$+$lower, no on-the-fly
claims.

\section{Metal-tick power analysis (parked)}

\frontierbox{Treat the live Metal FHE pipeline as the device under test.
Capture GPU-side power only while a tick-synchronized graph runs
\emph{controlled fixtures}. Classical DPA/CPA against aligned traces.
A distinguish or a clean negative are both science. Opportunistic capture
of arbitrary workloads is out of scope.}

First graduating experiment: fixture harness that emits tick markers and
a published ``no analysis yet'' log. Do not assign attacks against
third-party graphs.

\section{Rules for frontier essays}

\begin{enumerate}
  \item Label every paragraph frontier or first-experiment.
  \item Name the \claimkind{C} row it would become, and the \claimkind{N} it must not
        imply on the way.
  \item Cite Chapter~\ref{ch:living}. If the five cells cannot even be
        \emph{imagined}, the essay is not ready.
\end{enumerate}


\appendix
\part{Appendices}

\chapter{Notation}
\label{app:notation}

\begin{center}
\begin{tabular}{@{}lp{0.68\textwidth}@{}}
\toprule
Symbol & Meaning \\
\midrule
$q$ & Torus modulus; HELUT uses $q=2^{32}$. \\
$N$ & Polynomial degree / ciphertext shape. Not LUT count. \\
$n$ & LWE dimension (often $n=N$ in HELUT boolean production). \\
$B$ & Batch axis: independent copies of one netlist. \\
$\delta$ & $q/(2N)$, rotation / message spacing. \\
$\sigma$ & LWE Gaussian width. \\
$\varepsilon$ & Ingest failure probability (Gaussian certificate). \\
$B_{\mathrm{bk}}$ & BK noise bound; $0$ in HELUT now (\hid{4}). \\
$\ell$ & GGSW gadget levels. \\
$W$ & CMUX tile width in the Metal compiler. \\
$R_q$ & $\mathbb{Z}_q[X]/(X^N+1)$. \\
$M_A$ & Negacyclic Toeplitz matrix of polynomial $A$. \\
$w$ & TensorLUT INIT genome in $[0,1]^{64L}$. \\
$L$ & Number of LUT6 cells in a TensorLUT netlist. \\
$y(w)$ & Multilinear extension of INITs (exact on $\{0,1\}$). \\
$t$ & Target bit vector for $F_{\mathrm{crypto}}$. \\
$\pi(w)$ & Discreteness penalty $\sum_i w_i(1-w_i)$. \\
$E(w)$ & Emitter $\mathbf{1}[w\ge\tfrac12]$. \\
$M$ & Freeze mask: coordinates excluded from $\pi$. \\
$\lambda$ & Discreteness squeeze on $\pi(w)$; Theorem~1 needs $\lambda\ge 0$. \\
\texttt{\$lut} & Yosys look-up cell; one programmable gate. \\
SING & Encrypted $\equiv$ clear on generic stimuli. \\
\cid{$n$}, \hid{$n$}, \nid{$n$} & Claim, hedge, non-implication IDs. \\
\bottomrule
\end{tabular}
\end{center}

\chapter{Claim index (snapshot)}
\label{app:claims}

Canonical living inventory: \corpus{directives/claim-sheet.md}.
If this appendix disagrees with the sheet, the sheet wins.
Snapshot epoch: \livingepoch{}.

\section{Reproducible results}

\begin{center}
\small
\begin{tabular}{@{}lp{0.72\textwidth}@{}}
\toprule
ID & Result \\
\midrule
\cid{1} & Yosys \texttt{\$lut}/DFF $\to$ Metal/CPU tensor eval (oracle). \\
\cid{2} & Welchman path breaks known P1030684 end-to-end. \\
\cid{3} & $\le 10$-plug / SAT kill chain removes ghosts. \\
\cid{4} & Encrypted LWE/GLWE $+$ GGSW BR per \texttt{\$lut}. \\
\cid{5} & Certificates on encrypted ticks. \\
\cid{6} & Encrypted $\equiv$ clear multi-netlist SING; adder $N\le 1024$. \\
\cid{7} & Metal 1-LUT BR microbench (fused ancestor numbers). \\
\cid{8} & TensorLUT M4 baseline $F_{\mathrm{crypto}}=0$ $+$ Verilog emit. \\
\cid{9} & Stecker involution sandwich; blind 3-pair PASS. \\
\cid{10} & Enigma256 reciprocity; fail-closed / bijection. \\
\cid{11} & Potsdam/Plaice keys $\neq$ P1030680. \\
\cid{12} & Windowed discriminator vs whole-message turnover flaw. \\
\cid{13} & Metal tiled-kernel BR at $N=1024$ (pre-fusion lineage). \\
\cid{14} & Metal full\_adder SING at $N=1024$ (pre-NTT lineage). \\
\cid{15} & Metal netlist-scheduled SING at $N=1024$. \\
\cid{16} & Metal fused EP kernel; $1.043\,\mathrm{s}$/BR at $N=1024$. \\
\cid{17} & GPU-resident BR tile; $0.519\,\mathrm{s}$/BR; SING $12.2\,\mathrm{s}/8$. \\
\cid{18} & 3-prime NTT persist BR; $0.433\,\mathrm{s}$/BR; SING $15.6\,\mathrm{s}/8$. \\
\cid{19} & TensorLUT continuous$\to$discrete Theorem~1 (six machine-checked lemmas). \\
\cid{20} & Wavefront-parallel independent \texttt{\$lut} BRs; boolean SING $10.6\,\mathrm{s}/8$. \\
\bottomrule
\end{tabular}
\end{center}

\section{Open hedges}

\begin{center}
\small
\begin{tabular}{@{}lp{0.72\textwidth}@{}}
\toprule
ID & Asterisk \\
\midrule
\hid{1} & Calibrated bits $\neq$ lattice-estimator until Sage fills JSON. \\
\hid{2} & \textbf{Closed} 2026-08-12 ($Z_{2N}$ pack). Kept as history. \\
\hid{3} & \cid{20} wavefront-parallel NTT SING beats \cid{17}; crypto $\ell=2$ untimed; Sage \hid{1} pending. \\
\hid{4} & Noisy BK modeled; BK encrypt still noiseless. \\
\hid{5} & \texttt{*PublicMS} gadgets ($g_0=\delta$): on-lattice intent, not a closer of old \hid{2}. \\
\hid{6} & TensorLUT / quarantine vs campaign: parallel research, not P1030680 PT. \\
\hid{7} & Catalog / Regenbogen / UEBUNG: negatives graded; middle ring $\neq A$ untested. \\
\bottomrule
\end{tabular}
\end{center}

\section{Standing non-implications}

Do not let prose imply: new lattice assumption; ``we invented TFHE'';
production keys without estimator; mock-torus $=$ FHE; P1030680 decrypted;
campaign fitness $=$ encrypted tick rate; TensorLUT $=$ U-534 break;
side-channels measured; quantum attacks analyzed.
Wild ideas in Chapter~\ref{ch:frontier} are trajectory, not results.

\chapter{Fifteen-week syllabus}
\label{app:syllabus}

Suggested title: \emph{Reconfigurable Homomorphic Computing}.
3 credits. Mix of lecture, seminar, and lab (Chapter~\ref{ch:instructor}).
This table is a \emph{future} outline, not a schedule you can run.
The book is not ready to teach.

\begin{center}
\small
\begin{tabular}{@{}clp{0.58\textwidth}@{}}
\toprule
Wk & Part & Classroom \\
\midrule
1 & I & Subject definition; living boxes; claim audit assigned. \\
2 & II & Netlists, Yosys \texttt{\$lut}, $B$ vs $N$; Lab~1. \\
3 & II & Torus, LWE, PBS sketch; $N=4$ Toeplitz. \\
4 & III & Pillar~I pipeline; oracle vs encrypted; Lab~2. \\
5 & III & SING; $Z_{2N}$ pack as a worked bug; Lab~3. \\
6 & III & Metal compiler Phase~1/2; fused DNF as a case. \\
7 & III & Certificates; hardness table; Lab~4. \\
8 & --- & Midterm (foundations $+$ five-cell); Lab~5 optional. \\
9 & IV & TensorLUT setup; Theorem~1 (\cid{19}); Lab~6. \\
10 & IV & Shatter vs hold; involution sandwich; campaign as case study (no decrypt). \\
11 & V & E256 contract; Red/Blue; fail-closed. \\
12 & VI & Application gallery; student five-cell cards; C52--C54 vs remainders. \\
13 & VI & Labs catch-up; claim-audit drafts due. \\
14 & VI & Open hedges; pick-a-next-experiment. \\
15 & VI & Frontier seminars; essays due; living-epoch recap. \\
\bottomrule
\end{tabular}
\end{center}

Final: take-home claim audit plus a short oral defense of one reproduced
box and one hedge box.
If the student upgrades the hedge, they have not passed the course's
signature skill.

\chapter{Reproduce commands (course subset)}
\label{app:repro}

Canonical file: \corpus{REPRODUCE.md}.
Cookbook: \corpus{directives/parameter-cookbook.md}.
Build:

\begin{lstlisting}
swift build -c release --product helut
\end{lstlisting}

\section{Documents}

\begin{lstlisting}
make textbook   # this book
make paper      # three-pillar report
make writeup    # P1030680 campaign report
\end{lstlisting}

\section{Hardness (\hid{1})}

\begin{lstlisting}
.build/release/helut --hardness-table
.build/release/helut --estimator-export \
  > logs/helut-estimator-pending.json
python3 Scripts/helut_lattice_estimate.py \
  --pending logs/helut-estimator-pending.json \
  --out logs/helut-estimator-results.json
\end{lstlisting}

Until SageMath provides \texttt{sage.all}, results stay null.

\section{Encrypted SING (\cid{4}--\cid{6})}

\begin{lstlisting}
.build/release/helut --bench netlist.json --degree 8 \
  --bench-encrypted --cpu-only --sing --vectors 8
./Scripts/helut_encrypted_sing.sh
\end{lstlisting}

\section{Metal microbench and NTT (\cid{18})}

\begin{lstlisting}
make test-metal-p1
.build/release/helut --bench-encrypted-micro --degree 1024 \
  --trials 2 --warmup 1
\end{lstlisting}

Default at $N>64$ is tiled-kernel with inlined NTT.
Do not assign \texttt{--metal-br-fused} at $N=1024$ as a timed lab.

\section{TensorLUT Theorem 1 (\cid{19})}

\begin{lstlisting}
swift test -c release --filter testTensorLUTFormalCertificate
\end{lstlisting}

Six lemmas must \texttt{hold}. Statement:
\corpus{directives/tensorlut-theorem.md}.
Not a U-534 / P1030680 decrypt.

\section{Boolean oracle benches (\cid{1})}

\begin{lstlisting}
./Scripts/helut_boolean_bench.sh
./Scripts/helut_boolean_scale.sh
\end{lstlisting}


\backmatter
\bibliographystyle{plain}
\bibliography{refs}

\end{document}
